%% toc: P15 164 164.7,164.8
% Source-specific presentations: F. Eberhardt and C. Glymour, Hans Reichenbach's Probability Logic (HHL 10, ch. 10, pp. 357--389);
% I. Niiniluoto, The Development of the Hintikka Program (HHL 10, ch. 9, pp. 311--356).
\chapter{Two histories of inductive logic: Reichenbach's probability logic and the Hintikka program}
\label{sec:p15-w02b-p15-030}

This chapter reconstructs two chapters of the \emph{Inductive Logic}
volume of the \emph{Handbook of the History of Logic}. The first, by
Frederick Eberhardt and Clark Glymour, is a critical history of Hans
Reichenbach's frequency-based probability logic and of his attempt to
justify induction by the straight rule. The second, by Ilkka Niiniluoto,
follows Jaakko Hintikka's program of inductive logic from its Carnapian
background through the two-dimensional \(\lambda\)--\(\alpha\)
continuum and the \(K\)-dimensional system to its applications in
semantic information, acceptance, cognitive decision theory, theories,
analogy, truthlikeness and machine learning.

The purpose is twofold. A reader who meets either program in the
literature should be able to locate each definition, displayed formula,
argument and objection as these chapters present it, and to tell which
claims are the historical figure's, which are the chapter's own
reading, and which are mathematics that can be checked. The inputs are
the two chapters themselves. Their conclusions are of three kinds:
historical theses about the development of the programs; philosophical
verdicts (for example that Reichenbach's justification of the straight
rule is ``extremely weak'', or that the Hintikka program is progressive
in the sense of Lakatos); and formal results about probability
measures on monadic languages. Only the third kind is proved here.

Page locators refer to the camera-ready text of the two chapters
(printed copyright footer 2009, Elsevier BV). The Eberhardt--Glymour
copy carries printed pages 357--388; the final bibliography page 389 is
not in it. The Niiniluoto copy carries the complete pages 311--356.
Where the chapter quote or cite Reichenbach, Russell, Creary or
others, the page numbers given are the chapters' own locators for
those works, reported as they give them; the cited works themselves are
not reproduced here
\cite{EberhardtGlymour2011ReichenbachReading,Niiniluoto2011HintikkaProgramReading}.

Reichenbach's own 1937 French paper and the 1938 excerpt of
\emph{Experience and Prediction} are reconstructed separately in
\cref{sec:p15-reichenbach-probability-logic}; Carnap's 1950 system and
Zabell's history of it are in
\cref{sec:p15-carnap-logical-probability-and-induction}. Those chapters
are linked where they already prove a result that the present histories
use; nothing below replaces them.

\begin{hbconvention}{def:p15-w02b-p15-030-claim-types}{Kinds of claim in this chapter}
Four kinds of claim are kept apart.
\begin{enumerate}
\item \emph{Textual-historical}: what Reichenbach, Hintikka, Carnap and
  their critics wrote and when, as the two chapters report it, and the
  interpretation given in each chapter, which is marked as such.
\item \emph{Philosophical}: theses argued for or against, for example
  that probability is more fundamental than causality, that a blind
  posit needs no belief in its truth, or that inductive probabilities
  depend on extra-logical parameters.
\item \emph{Mathematical}: statements about probability measures,
  limits, constituents and parameters. Those stated as theorems below
  are proved in the annex under the hypotheses written in the statement.
  Formal claims that the chapters state without proof, and which are
  not proved here, are reported as the chapter's statement with its own
  attribution.
\item \emph{Methodological-empirical}: claims about scientific practice
  (for example that scientists rely on parametric assumptions, or that
  simulations are used to estimate convergence rates). The chapters give
  these as general observations, not as findings of a study with a
  stated design, sample or uncertainty; they are reported as such.
\end{enumerate}
\end{hbconvention}

% ---------------------------------------------------------------------
\section{Eberhardt and Glymour on Reichenbach's probability logic}
\label{sec:p15-w02b-p15-030-eg}

Eberhardt and Glymour set out to characterize Reichenbach's approach to
inductive reasoning through the influences that separate it from
neighbouring views, then to state the logic, its interpretation and its
justification, then to report the criticisms of Popper, Nagel and
Russell and Reichenbach's verdicts on rival programs, and finally to
give their own assessment. Their reconstruction rests mainly on the
German (1935) and English (1949) editions of \emph{The Theory of
Probability}, earlier \emph{Erkenntnis} papers, \emph{Experience and
Prediction}, and Reichenbach's replies to critics; they say that on some
points they could not determine what he intended
\cite[pp.~357, 387]{EberhardtGlymour2011ReichenbachReading}.

\subsection{Influences and the inverted approach}
\label{sec:p15-w02b-p15-030-eg-influences}

Eberhardt and Glymour call Reichenbach's account of induction thoroughly
empirical: an inductive logic is not to be built from logical
principles independent of experience but drawn out of the reasoning that
works in the natural sciences. They give two reasons for this
``inverse'' order. First, Heisenberg's results, including the
uncertainty principle, raised the question whether propositions fixing
both location and velocity have a determinate truth value; if residual
uncertainty is unavoidable there, the goals of induction may need to be
lowered, soundness may be an inappropriate demand, and binary truth
values may not apply. Second, the evidence for Einstein's relativity
showed that empirical findings can overturn claims once held necessary,
Euclidean geometry among them, despite Kant's transcendental arguments.
The foundations of any inference system therefore had to be shown true
of the world rather than assumed a priori. Reichenbach called the
confirmation of the fit between formal structures and the world
``coordination''; Eberhardt and Glymour remark that ``alignment'' would have
conveyed his meaning better \cite[pp.~357--358]{EberhardtGlymour2011ReichenbachReading}.

Causality is the case in which these influences acted most strongly.
In the 1915 doctoral thesis, under Kantian influence, causal knowledge
counted as synthetic a priori. With the a priori view of space gone,
Reichenbach gave up the a priori status of causality as well; causal
knowledge then had to be acquired empirically, which required an
inductive procedure, and that procedure had to handle the uncertainty of
empirical data (measurement error, illusion). Heisenberg's results
added a further problem: microphysical uncertainty seems permanent while
macrophysics uses stable causal relations. Eberhardt and Glymour conclude that
causal knowledge remained Reichenbach's paradigm case for induction and
that probability sat at its base
\cite[p.~358]{EberhardtGlymour2011ReichenbachReading}.

The same pressures changed his view of deductive logic. Although he did
not treat the foundations of logic in detail, he rejected the view that
logic records necessary truths about things: necessity comes from the
syntactic rules of a language, and reality need not conform to them (Eberhardt and Glymour cite Reichenbach 1948). He adopted instead a formalist,
Hilbert-style view: an inference system is an axiom system, its theorems
are what valid deductions from the axioms yield, and whether they are
true of the world depends on how well the axioms can be coordinated with
it, which is an empirical matter of success or usefulness. In the light
of quantum theory he rejected classical logic altogether (they cite
Reichenbach 1944). Applied to induction, this gives an axiomatic method:
start from causal knowledge we do not doubt, study the scientific
methods by which it was gained, and make explicit, as axioms, the
assumptions behind the cases in which Hume's inductive gap seems to have
been partly bridged, together with the stories we tell about why those
assumptions are reliable \cite[pp.~358--359]{EberhardtGlymour2011ReichenbachReading}.

Two influences of a second tier follow. The logical positivists matter,
but Eberhardt and Glymour call Reichenbach's position distinctly ``negativist'':
he denied certainty even for primitive perception, yet, unlike Popper,
held that once uncertainty is taken into account a scientific hypothesis
can be given a positive probability. Richard von Mises, Reichenbach's
colleague in Berlin and Istanbul, is named the largest influence on his
concept of probability; from this came the limiting-frequency account
and the straight rule, which Eberhardt and Glymour regard as his best-known and
most disputed legacy. They add that piecing together a framework
developed over a lifetime is hard, that Reichenbach did not always write
clearly and sometimes revised his views in ways that did not make them
clearer, and that the revised points are of particular interest because
they show what he found hardest to settle and what he thought still
needed work. Sitting, in their image, between the thrones of the
established philosophy of his day, he anticipated many later ideas, which
is a further reason for interest in his views
\cite[pp.~359--360]{EberhardtGlymour2011ReichenbachReading}.

\subsection{The basic set-up}
\label{sec:p15-w02b-p15-030-eg-setup}

On Eberhardt and Glymour's account Reichenbach separates deductive and mathematical
logic, which concern relations among tautologies and rest on an axiom
system whose truth of the world is a secondary question, from inductive
logic, which concerns truth in reality. That we seem unable to think
except by certain inferences no more makes them true of the world than
our inability to picture non-Euclidean space makes real space
Euclidean. Inductive logic has to establish whether the mathematical
relations stated by a law hold between the real quantities its symbols
stand for. Deductive semantics can be adjusted to fit syntax; the world
allows no such freedom. Inductive logic therefore needs both a formal
semantics and a mapping between the syntactic law and the quantities it
is about \cite[p.~360]{EberhardtGlymour2011ReichenbachReading}.

Reichenbach is a realist, but access to the world is indirect: he
compares us to observers inside a cloth cube who infer the objects
outside from the shadows cast on its walls (Eberhardt and Glymour cite
\emph{Experience and Prediction}, 1938). Access is also inexact. Data are
``unclean'' for two reasons: countless small influences add noise to any
measurement, and, if Heisenberg's principle expresses a real rather than
an epistemic indeterminacy (a reading Eberhardt and Glymour attribute to
Reichenbach despite his friendship with Einstein), exact measurement is
impossible in principle. Hence all laws rest on probabilistic inference
from samples, and an inductive logic cannot be two-valued: its values
lie between 0 and 1, the ``truth value'' being the probabilistic
support. This is what Reichenbach calls probability logic
\cite[pp.~360--361]{EberhardtGlymour2011ReichenbachReading}.

Support is again modelled on scientific method. A law is a universal
claim based on a finite run of observations, and its uncertainty comes
at least partly from not knowing whether later observations will follow
the pattern so far. So the probability of an event is the limit of the
relative frequency of events of the same type in an infinite sequence;
what ``the same type'' means becomes one of Reichenbach's central
problems, and Eberhardt and Glymour set aside some further conditions on the
sequence at this point \cite[p.~361]{EberhardtGlymour2011ReichenbachReading}.

\begin{hbdefinition}{def:p15-w02b-p15-030-eg-coordination}{Coordination of events and propositions (as Eberhardt and Glymour present it)}
Let \(E_1,E_2,\ldots\) be a sequence of events and let \(S_i\) be the
proposition describing \(E_i\). Reichenbach's claim, as reported, is
that the sequence of propositions reproduces the structure of the
sequence of events isomorphically, so that a probability value can be
attached interchangeably to an event type or to the propositions
describing it. The probability of a proposition is the limiting
frequency of its occurrence in a sequence of propositions coordinated
with a sequence of events in the world. For a ball rolled repeatedly
down an inclined plane, with each descent timed, proposition~1 may
state that the ball took 4.2 seconds \((\pm\delta)\); its probability
is the limiting frequency of that proposition among the propositions
describing the trials. Because continuous quantities can only be
reported within error intervals, such a limiting frequency can be
nonzero; Eberhardt and Glymour add that Reichenbach is vague on this point and
describe the construction as a formal version of building a histogram.
Its consequence is that the probability logic must be a logic of
sequences \cite[pp.~361--362]{EberhardtGlymour2011ReichenbachReading}.
\end{hbdefinition}

Empirical sequences are finite, so the limit is left open, and any
estimation method needs a justification over its rivals. Reichenbach's
estimate is simple: treat the frequencies in the available initial
segment as if they were those of the limit, and revise them as the
sequence grows. This is the straight rule. Its justification, as Eberhardt and Glymour reconstruct it, has three parts.

\begin{hbdefinition}{def:p15-w02b-p15-030-eg-three-part}{The three-part justification of the straight rule}
\begin{enumerate}
\item \emph{Higher-order probabilities.} Probabilities of a higher
  order, supposedly drawn from more general and abstract knowledge, give
  reason to believe that the empirical distribution is approximately
  accurate.
\item \emph{Truncation by blind posits.} A hierarchy of higher-order
  probabilities, none of them certain, need not produce an infinite
  regress of values we cannot determine: at some level the regress is
  cut off by blind posits, wagers used in place of probability values.
\item \emph{Convergence.} The straight-rule estimate converges to the
  limiting frequency to any desired accuracy after a finite amount of
  data. Reichenbach grants that other rules have the same guarantee but
  calls the straight rule ``simpler''.
\end{enumerate}
\cite[p.~362]{EberhardtGlymour2011ReichenbachReading}
\end{hbdefinition}

If these claims are granted, the result is a logic of probabilistic
inference that relates sequences of propositions rather than single
propositions. It gives each sequence a continuous value, the limit of
its relative frequencies; for empirical claims this limit is estimated
by evidence and the straight rule, and two-valued logic is recovered as
a limiting case. Reichenbach regarded this as the most that Hume's gap
permits: certainty is replaced by high probability, the continuous
values express graded truth about the world, and the result vindicates
the rationality of scientific inference, since it licenses more
confidence in some claims than in others according to sample size or to
similar findings in related fields. The probability calculus supplies
the machinery that transfers support between claims; hoping for a
stronger bridge over the gap is, in his view, wishful thinking
\cite[p.~362]{EberhardtGlymour2011ReichenbachReading}.

\subsection{Probabilities as limiting frequencies: the chronology}
\label{sec:p15-w02b-p15-030-eg-frequency}

For Reichenbach scientific probability is an objective quantity, not a
degree of belief. The difficulty is to say exactly what that quantity
is while giving grounds for probability judgments, and his answer
changed over his life \cite[p.~363]{EberhardtGlymour2011ReichenbachReading}.

\paragraph{The 1915 thesis.} In the doctoral thesis the probability of
an event is its relative frequency in an infinite sequence of causally
independent and causally identical trials. Under the influence of the
neo-Kantians (Cassirer, Natorp) causality is a primitive concept, prior
to probability, and causal knowledge of individual events is synthetic a
priori, supposedly established by Kant's transcendental deduction of the
causal principle. That knowledge fixes causal independence and identity
of circumstances, and so offers a non-circular, objective basis for
probability. If causally independent identical trials were shown to be
probabilistically independent and identically distributed, the law of
large numbers would give convergence of the empirical distribution to
the true one \emph{in probability}. Eberhardt and Glymour's footnote states that law:
in independent identical trials, for every \(\varepsilon>0\) the
probability that the observed success frequency differs from the true
success probability by more than \(\varepsilon\) tends to zero as the
number \(n\) of trials grows. Reichenbach knew this weak law but judged
convergence in probability too weak, because it would put probability
into the definiens of probability and make the definition circular; he
wanted convergence with certainty
\cite[p.~363 and n.~1]{EberhardtGlymour2011ReichenbachReading}.

To get it, the 1915 thesis gives a transcendental argument for a
synthetic a priori \emph{principle of lawful distribution} guaranteeing
that every empirical distribution converges. The argument: without such
a principle the knowledge recorded in laws of nature would be
impossible; science plainly has knowledge of lawful causal relations;
laws state general regularities while Kantian causal knowledge concerns
single events, so something must aggregate the single-event knowledge
into laws; hence the principle exists. Given it, convergence is certain,
and even without knowing when or how fast it occurs we are on the right
track in using the empirical distribution
\cite[p.~363]{EberhardtGlymour2011ReichenbachReading}.

Eberhardt and Glymour find this unconvincing even with its gaps filled (for
example the step from causal to probabilistic independence). Granting
the now implausible view that causal knowledge of token events is
synthetic a priori, it is simply false that there is a synthetic a
priori guarantee that sequences converge to a limiting distribution:
many sequences have relative frequencies without a limit
\cite[p.~364]{EberhardtGlymour2011ReichenbachReading}. (Concrete binary
sequences whose frequencies have no limit occur in item~1 of
\cref{thm:p15-reichenbach-frequency-boundaries}.)

\paragraph{The 1927 notes.} The argument changed substantially between
1915 and the English \emph{Theory of Probability} of 1949 (Eberhardt and Glymour
refer to the introduction of the 2008 translation of the thesis for the
details). In notes of 1927, catalogued as HR 044-06-21 in the
Reichenbach Collection, and referring back to discussions with Paul
Hertz, Reichenbach concluded that convergence with certainty cannot be
had and that only convergence in probability can be guaranteed. Eberhardt and Glymour call this a transitional thought: it is an observation about the
law of large numbers, which presupposes a distribution from which
initial segments are drawn independently, whereas the limiting-frequency
interpretation gives the straight rule certain convergence provided a
limit exists. He also reversed the order of primitives. Once Einstein
had undermined the a priori status of space and time, Reichenbach
reconsidered that of causality and concluded that probability is more
basic: causal relations are inferred from probabilistic relations plus
further assumptions. Single-event (``actual'') causal claims became
elliptical, either tacitly about a sequence or a fictitious transfer of
a type-level claim to a token
\cite[p.~364 and nn.~2--3]{EberhardtGlymour2011ReichenbachReading}.

\paragraph{Von Mises, randomness and normal sequences.} Without causal
primitives a new foundation for probability was needed, and here
Reichenbach's concern met von Mises's attempt (1919) to found
probability on random sequences. Both wanted to reduce probability to a
property of infinite sequences and so avoid circularity. Not every
sequence is admissible: in the alternation \(1,0,1,0,\ldots\) the
limiting frequency of 1 is \(1/2\), yet after a long enough segment the
next term is known for certain. Von Mises therefore restricted attention
to random sequences, in which no betting system on the next term given
the earlier ones makes money. Two conditions were taken as necessary:
lack of after-effect (the probability of an outcome does not depend on
the preceding outcomes) and invariance under subsequence selection by
any rule that depends only on the indices. Writers on randomness spell these out
differently \cite[pp.~364--365]{EberhardtGlymour2011ReichenbachReading}.

Reichenbach rejected random sequences because he saw no prospect of
capturing randomness formally; there were known difficulties in showing
all the conditions jointly satisfiable, some of which he had pointed out
(1932). (In a footnote Eberhardt and Glymour record that Nagel, in 1936, agreed
about this but thought von Mises's weaker notion could still be
formalized.) Reichenbach settled for \emph{normal sequences}, a strictly
larger class than the random ones.

\begin{hbdefinition}{def:p15-w02b-p15-030-eg-normal}{Normal sequence (as reported)}
A sequence of events is \emph{normal} if it is free of after-effect and
the probabilities of its event types are invariant under regular
divisions. In a sequence with after-effect, roughly, an event \(E\) at
index \(i\) gives later indices probabilities different from the
limiting relative frequencies of the events there; Eberhardt and Glymour say
Reichenbach's definition of after-effect is not entirely clear. A
regular division selects every \(k\)-th term of the sequence for some
fixed \(k\); Eberhardt and Glymour note that the actual conditions are more
intricate. The probability of \(E\) is then the limiting relative
frequency of \(E\) in a normal sequence
\cite[p.~365 and n.~4]{EberhardtGlymour2011ReichenbachReading}.
\end{hbdefinition}

The original French presentation of normality, phase division and the
missing-limit cases is reconstructed in
\cref{sec:p15-reichenbach-normality}.

\paragraph{Finite segments and higher-order probability.} As an abstract
definition this works, but it does not settle scientific probabilities,
because measurement sequences are finite and a finite initial segment
carries no information about the limit. Reichenbach nevertheless tells
us to treat the empirical distribution as roughly equal to the limiting
one, appealing to a higher-order probability that the limit lies within
a narrow band of width \(\delta\) around the observed frequency. That
higher-order probability is itself empirical but indirect: it comes from
a sequence of first-order values, that is, from a sequence of sequences,
and so integrates the results of different inductions
\cite[p.~365]{EberhardtGlymour2011ReichenbachReading}.

\begin{hbexample}{def:p15-w02b-p15-030-eg-gravitation}{Cross-induction for the gravitational constant}
As Eberhardt and Glymour render Reichenbach's example (which they locate
in the 1949 edition, pp.~438--440), it runs as follows. Measurements
\(M_1,M_2,\ldots,M_n\) of the gravitational constant are classified by
the indicator
\[
 X=I(|M-\gamma|\leq\delta),
\]
which is 1 when a measurement falls in the band \(\gamma\pm\delta\) and 0
otherwise, giving a 0--1 sequence such as \(1,1,0,1,0,\ldots,0\). If the
sequence were infinite it would have limiting distribution
\(P(X=1)=p\), \(P(X=0)=1-p\); the available measurements give the
empirical distribution \(\hat P(X=1)=\hat p\), \(\hat P(X=0)=1-\hat p\).
Reichenbach asserts a higher-order probability \(q\) with
\(P(|\hat p-p|<\varepsilon)=q\) for small \(\varepsilon\), and claims
that \(q\) can be estimated from several measurement sequences, for
example (a) on the moon, (b) on some planet, and (c) with the Cavendish
balance:
\[
\begin{array}{ll}
 a: & 1,1,1,0,0,\ldots,1\\
 b: & 0,1,1,1,1,\ldots,0\\
 c: & 0,0,1,0,0,\ldots,1
\end{array}
\]
Their empirical frequencies \(\hat p_a,\hat p_b,\hat p_c\) form a new
sequence. With \(\hat p_a=0.8\), \(\hat p_b=0.7\), \(\hat p_c=0.79\) and the
band \(0.8\pm\varepsilon\), \(\varepsilon=0.05\), two of the three values
fall inside, so \(\hat q=2/3\). The frequency of 1s in one row is the
probability that the statement about the gravitational constant is true
for that test object; the second-order \(q\), taken across rows, is the
probability that the first-order probability statement is true. This
mutual validation across sequences is what is supposed to supply a
probability of convergence for any empirical distribution
\cite[pp.~365--367]{EberhardtGlymour2011ReichenbachReading}. The
arithmetic, and its limits, are discussed in
\cref{sec:p15-reichenbach-cross-inductions}.
\end{hbexample}

In a parallel example about the melting points of metals, the fact that
many metals melt is a reason to expect that a metal not yet seen to melt
has a melting point. Its own measurements point to probability zero, but
the second-order probability that its empirical distribution indicates
its limit is very low, because that second-order value integrates the
other metals. Reichenbach calls such cross-inductions a ``network of
inductions'' \cite[p.~367]{EberhardtGlymour2011ReichenbachReading}.

\paragraph{The hierarchical Bayesian comparison.} Eberhardt and Glymour offer a
second way to view the approach. In a hierarchical Bayesian procedure,
measurements are used to estimate parameters of the distribution of a
quantity; those parameters are themselves given a higher-order
distribution with hyper-parameters; several measurement sequences
estimate the hyper-parameters; and the lower-level parameters are then
recomputed given those estimates. Information flows between sequences
through the hyper-parameters, and the construction can be continued to
hyper-hyper-parameters, until one asks whether there are enough data to
estimate the top level. At that point Reichenbach puts blind posits in
place of probability estimates to stop the regress. Hierarchical
methods, Eberhardt and Glymour say, work well in theory but depend on deciding
whether and in what sense events are similar enough to be pooled in one
sample; for Reichenbach the corresponding question is the choice of
reference class \cite[p.~367]{EberhardtGlymour2011ReichenbachReading}.
The comparison is Eberhardt and Glymour's interpretive device; it is not
Reichenbach's own formalism.

\paragraph{Reference classes and the probability of a theory.}
Reichenbach goes so far as to assign probabilities to theories. To find
the probability that Newton's law of gravitation holds universally
(rather than for one test object), all measurements of the gravitational
constant go into one row, and the other rows are filled with
observations bearing on other laws, such as conservation of energy and
the law of entropy (Eberhardt and Glymour quote the 1949 edition, p.~439f). His
defence (p.~440 there) is that this reference class mirrors how
theories are actually judged: success of other laws raises confidence in
a given law, and failures of some laws are a reason to restrict others
not yet refuted. His illustration, in the words Eberhardt and Glymour quote:
\enquote{the fact that Maxwell's equations do not apply to Bohr's atom
is regarded as a reason to question the applicability of Newton's or
Einstein's law of gravitation to the quantum domain}
\cite[pp.~367--368]{EberhardtGlymour2011ReichenbachReading}.

Eberhardt and Glymour reply that the choice of reference class looks arbitrary,
and that they cannot see why the failure of Maxwell's equations for a
model of the atom should tell against Newton's equations, which are
themselves incompatible with Maxwell's, nor what reference class could
carry such a transfer. More generally Reichenbach offers no criterion:
events should be of the same type but not so alike that the variation
of interest disappears. He recommends the narrowest class with stable
statistics, stability being judged at the level of ``advanced
knowledge''. In a footnote Eberhardt and Glymour report his own illustration: the
measurements for Mercury converge to a different limit from those for
the other planets, which should lead to splitting the class of planets
into those near the sun and those farther away (1949 edition, p.~439).
Eberhardt and Glymour object that the recommendation begs the question, since
stable limits are what is to be found, and that trivially stable
statistics are always available in the narrowest nonempty class, a
single event, which cannot be what Reichenbach meant. Reading the appeal
to advanced knowledge as an appeal to blind posits invites the question
why one should not guess the lower-level frequencies directly. They
leave open whether an intuitive reference class could be formalized and
would then yield objective probabilities
\cite[p.~368 and n.~5]{EberhardtGlymour2011ReichenbachReading}.

Their summary: probability is a property of infinite normal sequences,
which share many features of random sequences; inferences are built on
finite initial segments; we are on the right track as long as
higher-order probabilities look promising, which they do when inductions
from many different measurement sequences agree; probabilities of every
order are objective because they are relative frequencies. The account
founders on the choice of reference class. The 1915 thesis had an easy
answer, the class of causally independent and identical trials backed by
a priori knowledge; having abandoned the a priori shore, in Eberhardt and Glymour's
image, Reichenbach was at sea
\cite[pp.~368--369]{EberhardtGlymour2011ReichenbachReading}.

\subsection{The probability logic: syntax, semantics and the rule of induction}
\label{sec:p15-w02b-p15-030-eg-logic}

With continuous values in place of truth values, Reichenbach claims that
every form of uncertainty in scientific inference can be represented in
a formal framework. Axiomatization places the logic in Hilbert's
formalist tradition: it needs only the axioms of propositional logic and
an interpretation of probabilities as properties of infinite sequences,
so that inductive reasoning reduces to deduction plus induction by
enumeration on suitable sequences. Tautologies come out as a special
case, because the sequences relevant to deductive logic are constant and
their properties can be determined with certainty; empirical sequences
are given only extensionally, so their values lie strictly between 0
and 1. The inferences follow the probability calculus, and their
application is justified by a convergence argument. A footnote asks what
probability, if any, Reichenbach would have given to a mathematical
statement not yet decided, such as Goldbach's conjecture
\cite[p.~369 and n.~6]{EberhardtGlymour2011ReichenbachReading}.

The syntax is that of propositional logic with conjunction,
disjunction, negation, implication and equivalence, but a proposition
receives a degree of probability \(p\) rather than truth or falsity.
Since the aim is to capture scientific probability, \(p\) is to be read
as a limit of relative frequency. Unlike propositional logic the system
cannot be compositional: the probabilities of \(A\) and of \(B\) do not
fix that of \(A\lor B\), just as marginals do not fix a joint
distribution. A third quantity is needed, and Reichenbach takes the
conditional probability of \(B\) given \(A\), because it can be read as
a subsequence selection. In a footnote Eberhardt and Glymour note that Russell,
Tarski and Nagel call a compositional logic ``extensional'' and its
opposite ``intensional''; they find the terminology misleading because
``extensional'' is also used of sequences given only by enumeration, and
they report Reichenbach's answer that his logic fits neither side
\cite[p.~370 and n.~7]{EberhardtGlymour2011ReichenbachReading}.

\begin{hbdefinition}{def:p15-w02b-p15-030-eg-composite}{The composite formulas as printed}
Given the marginals and the conditional, the value of any composite
formula built with the binary connectives is fixed by the probability
calculus. Eberhardt and Glymour's examples, as printed, are
\[
\begin{aligned}
 P(A\lor B)&=P(A)+P(B)-P(A)(B\mid A),\\
 P(A\equiv B)&=1-P(A)-P(B)+2P(A)P(B\mid A).
\end{aligned}
\]
In the first line the factor printed \((B\mid A)\) is to be read as
\(P(B\mid A)\), as in the second line
\cite[p.~370]{EberhardtGlymour2011ReichenbachReading}.
\end{hbdefinition}

With \(u=P(A)\), \(v=P(B)\) and \(w=P(A)P(B\mid A)=P(A\land B)\) for
\(u>0\), the two printed formulas are the disjunction and equivalence
rows \(u+v-w\) and \(1-u-v+2w\) of the three-parameter table proved in
\cref{thm:p15-reichenbach-coupling}. The same theorem gives the exact
range \(\max(0,u+v-1)\leq w\leq\min(u,v)\) of the coupling that the
marginals leave free, which is the precise content of the
non-compositionality Eberhardt and Glymour describe.

The propositional axioms are supplemented by axioms that Reichenbach had
developed for the mathematical probability calculus (Eberhardt and Glymour cite
the 1949 edition, pp.~53--65). In his notation
\(A\mathrel{\underset{p}{\supset}}B\) is the probability implication
``\(A\) implies \(B\) with degree \(p\)''; the dot is conjunction, the bar
over a class is complement, and \((\bar A)\) in parentheses expresses
that the class \(A\) is empty.

\begin{hbdefinition}{def:p15-w02b-p15-030-eg-axioms}{Reichenbach's probability axioms in Eberhardt and Glymour's display}
\begin{description}
\item[Univocality.]
\(\displaystyle (p\neq q)\supset\bigl[(A\mathrel{\underset{p}{\supset}}B).(A\mathrel{\underset{q}{\supset}}B)\equiv(\bar A)\bigr]\)
\item[Normalization.]
\(\displaystyle (A\supset B)\supset(\exists p)(A\mathrel{\underset{p}{\supset}}B).(p=1)\)\\
\(\displaystyle \overline{(\bar A)}.(A\mathrel{\underset{p}{\supset}}B)\supset(p\geq0)\)
\item[Addition.]
\(\displaystyle (A\mathrel{\underset{p}{\supset}}B).(A\mathrel{\underset{q}{\supset}}C).(A.B\supset\bar C)\supset(\exists r)(A\mathrel{\underset{r}{\supset}}B\lor C).(r=p+q)\)
\item[Multiplication.]
\(\displaystyle (A\mathrel{\underset{p}{\supset}}B).(A.B\mathrel{\underset{u}{\supset}}C)\supset(\exists w)(A\mathrel{\underset{w}{\supset}}B.C).(w=p\cdot u)\)
\end{description}
Eberhardt and Glymour gloss them as follows. The first makes the value of a
probability unique (two different values can hold together only when the
reference class is empty). The second gives every probability with a
nonempty conditioning class a value between 0 and 1 inclusive (a formal
implication gives degree 1; nonempty reference gives a nonnegative
degree). The third makes the probability of mutually exclusive events
the sum of their probabilities. The fourth is the chain rule
\(P(C,B\mid A)=P(C\mid B,A)P(B\mid A)\)
\cite[pp.~370--371]{EberhardtGlymour2011ReichenbachReading}.
\end{hbdefinition}

The existential clauses \((\exists r)\) and \((\exists w)\) assert that a
degree exists and has the stated value; they quantify over real numbers,
not over sequences. The 1937 French schemata, whose surface differs from
this display, are reconstructed in
\cref{def:p15-reichenbach-probability-implication}.

Eberhardt and Glymour compare the first three axioms with Kolmogorov's 1933
axioms, with one difference: the addition axiom secures only
\emph{finite} additivity. In Kolmogorov's system the chain rule follows
from the other axioms, whereas Reichenbach needs a separate fourth axiom
to pass between logical conjunction and arithmetical multiplication.
They add that the axioms do not yet axiomatize probability, because
nothing ensures that the domain is closed under complement and countable
union, that is, forms a sigma-field, and they report that van Fraassen
(1979) showed limiting relative frequencies in infinite sequences do not
satisfy these closure conditions \cite[p.~371]{EberhardtGlymour2011ReichenbachReading}.
An elementary instance of the phenomenon, two binary sequences with
frequency limit \(1/2\) whose conjunction and disjunction have no
frequency limit, is item~1 of
\cref{thm:p15-reichenbach-frequency-boundaries}; it is not van Fraassen's
own construction, which the chapter does not reproduce.

\begin{hbdefinition}{def:p15-w02b-p15-030-eg-rule}{The rule of induction}
The axioms are extended by one rule, the rule of induction or straight
rule, which Eberhardt and Glymour quote from the 1949 edition, p.~446:
\enquote{If an initial section of \(n\) elements of a sequence \(x_i\) is
given, resulting in the frequency \(f^n\), and if, furthermore, nothing is
known about the probability of the second level for the occurrence of a
certain limit \(p\), we posit that the frequency \(f^i\) \((i>n)\) will
approach a limit \(p\) within \(f^n\pm\delta\) when the sequence is
continued.}
\cite[p.~371]{EberhardtGlymour2011ReichenbachReading}
\end{hbdefinition}

Reichenbach takes this rule to be the only addition the formal logic
needs to produce inductive inferences: by the axioms, every inductive
inference about complex claims reduces to this simple enumeration. A
footnote records that he took the rule to be an instance of
C.~S.~Peirce's self-correcting method. The rule belongs to the
metalanguage and cannot, like other deductive rules, be reduced to the
object language, so it does not share the logical form of the axioms
\cite[pp.~371--372 and n.~8]{EberhardtGlymour2011ReichenbachReading}.

Eberhardt and Glymour then point out what the surface hides. If the axioms built a
logic from sequences of propositions and carried out the probability
calculus as subsequence selection on infinite sequences, Reichenbach
would have a purely syntactic inductive calculus covering probabilistic
and deductive inference alike, based on enumeration. But the axioms and
the rule mix notations that are never fully specified. The probability
axioms use arithmetic relations and existential quantifiers over real
variables, so arithmetic must belong to the language; yet neither that
machinery nor its first-order part is extended to the variables standing
for sequences, which occur in open formulas, presumably universally
quantified, in a propositional language enlarged by probability
implication. Reichenbach never states the joint syntax of the two systems. The semantics is
non-standard: a proposition's value is the limiting frequency in its
sequence, and a complex formula's value is computed from the frequencies
of the component sequences and of the subsequence for the conditional.
This only hints at a formal semantics. A formula with iterated
probability implications such as
\[
 (A\mathrel{\underset{p}{\supset}}B)\mathrel{\underset{q}{\supset}}C
\]
must either be excluded, which seems contrary to Reichenbach's aim of
covering all probabilistic inference, or be read by subsequence
selection, which is ill-defined because the antecedent
\(A\mathrel{\underset{p}{\supset}}B\) is not obviously the kind of object
that has a sequence interpretation \cite[p.~372]{EberhardtGlymour2011ReichenbachReading}.

\subsection{Interpretation: posits and single cases}
\label{sec:p15-w02b-p15-030-eg-interpretation}

Many-valued and continuously valued logics for modality were proposed in
the same decades; Reichenbach thought them largely misguided because
they became formal constructions with little relation to modal talk in
natural language (Eberhardt and Glymour cite his 1934 paper). By contrast he saw
his logic, through its inclusion of ordinary probabilistic reasoning, as
a rational reconstruction of scientific inference. Its semantics was, he
thought, already ``coordinated'' with science, since the foundation of
probability supplied the passage from evidence to complex propositions
and back: data are a sequence of events, scientific probabilities are
relative frequencies in it, each proposition describes one datum, and
the resulting sequence of propositions is isomorphic to the sequence of
events \cite[pp.~372--373]{EberhardtGlymour2011ReichenbachReading}.

Two gaps remain. The logic is a calculus of infinite sequences, but
scientific data are finite; and natural language, and perhaps science,
gives probabilities to singular propositions that have no obvious
sequence. For singular propositions Reichenbach's answer is the posit, a
wager (Eberhardt and Glymour say: not in de Finetti's strict sense).

\begin{hbdefinition}{def:p15-w02b-p15-030-eg-posits}{Blind and appraised posits (as reported)}
A posit attached to a single proposition is \emph{blind} when no data
bear on its probability and \emph{appraised} once evidence is available,
in which case it should equal the relative frequency of the relevant
event in the data. Reichenbach gives no explicit constraint on the form
of a blind posit; Eberhardt and Glymour judge that implicitly it should resemble a
flat prior, giving every possibility equal probability
\cite[p.~373]{EberhardtGlymour2011ReichenbachReading}.
\end{hbdefinition}

Reichenbach's example (1949 edition, p.~366f, as Eberhardt and Glymour cite it) is
the proposition that Caesar was in Britain at a given time. With no
evidence its probability \(p\) is a blind posit. One relevant sequence is
the sequence of historians' reports of Caesar's movements; the
probability becomes the frequency with which his presence in Britain at
that time is reported, and the posit becomes appraised. Single-case
probabilities are thus fictive or elliptical, referring tacitly to
sequences. Eberhardt and Glymour find this unsatisfying; in a footnote they ask
how Reichenbach would distinguish silence about Caesar's whereabouts
from an explicit report of absence or presence, and say that reference
class selection might be meant to handle this but that no procedure is
given. They think a scientific case fits better: the probability that a
given atom decays after time \(t\) can more reasonably be read as the
frequency of decay after \(t\) in a set of atoms, a larger population
even if not a sequence \cite[pp.~373--374 and n.~9]{EberhardtGlymour2011ReichenbachReading}.

Even so, the reference class is finite, which returns us to the first
gap. Reichenbach's answer is the rule of induction: posit the observed
frequency of a finite sequence as its limit. This is a blind posit, since
nothing says the empirical distribution indicates the limit; it becomes
appraised when there are several sequences of a similar type. Science
and the logic alike proceed \emph{as if} the finite segment were
infinite, and aggregating data into different sequences is supposed to
turn blind posits into appraised ones, ``as if many blind eyes make
vision'' in Eberhardt and Glymour's phrase. Reichenbach argues, they report (citing
the 1949 edition, pp.~326--333 and~441), that under a flat prior his rule
of induction agrees with a Bayesian update, at least for the first
update, and that appraised posits are simply informative priors; his
hierarchy then has the structure of hierarchical Bayesian methods, which
did not yet exist \cite[p.~374 and n.~10]{EberhardtGlymour2011ReichenbachReading}.

\begin{hbtheorem}{thm:p15-w02b-p15-030-flat-prior}{proved=annex:p15-w02b-p15-030-flat-prior}{What a flat prior does and does not reproduce}
Let \(\theta\in[0,1]\) have the uniform prior density and let \(n\geq1\)
conditionally independent Bernoulli\((\theta)\) trials give \(k\)
successes. Then the posterior density is proportional to
\(\theta^k(1-\theta)^{n-k}\) on \([0,1]\), and its unique maximizer (the
posterior mode) is the straight-rule value \(k/n\). The posterior mean
\((k+1)/(n+2)\) equals \(k/n\) only when \(k/n=1/2\). In particular after
a single success the mode is \(1\) and the mean is \(2/3\).
\end{hbtheorem}

The theorem states one precise sense in which the claimed agreement
holds (the most probable value after one flat-prior update is the
observed frequency) and one in which it fails (the expected value is
Laplace's rule, not the frequency). The posterior mean itself is the case
\(\alpha=\beta=1\) of the beta updating in
\cref{thm:p15-polya-posterior-prediction}; the related Carnapian
calculation for \(c^*\) appears in \cref{thm:p15-carnap-predictive}. Eberhardt and Glymour's
own reservations are different and apply even under a flat prior at the
lowest level. How a posit becomes appraised, and why its appraised value
is unique, is unclear. If measurement sequences can be arranged in ways
that change which events count as first-order (they point to the
gravitational-constant example as one such rearrangement), the rule of
induction, which reflects empirical frequencies, will disagree on
higher-order claims with Bayesian updating. It seems the rule should be
used only at the most basic data level, preserving the objectivity of
probability there, while higher-order claims about convergence and the
integration of ``advanced knowledge'' proceed by Bayesian updating.
Reichenbach never discusses this mixture of updating methods
\cite[p.~374 and n.~11]{EberhardtGlymour2011ReichenbachReading}.

\subsection{Justification}
\label{sec:p15-w02b-p15-030-eg-justification}

The convergence argument, presumably his most disputed legacy, says that
if we use the empirical distribution as if it were the limit and keep
revising it as data arrive, our probability judgments converge to the
true probabilities provided the empirical distribution has a limit. The
proviso is essential, because not every infinite sequence has limiting
relative frequencies. Given a limit, for every \(\varepsilon\) some \(N\)
makes the empirical distribution within \(\varepsilon\) of the limit
thereafter; but since only the existence of a limit is assumed, and even
that conditionally, what is guaranteed is \emph{pointwise} convergence
(for each \(\varepsilon\) and each limiting frequency some \(N\)), not
\emph{uniform} convergence (for each \(\varepsilon\) one \(N\) for all
distributions). Uniform convergence permits confidence intervals and
rates; pointwise convergence does not. Reichenbach's \emph{principle of
finite attainability} (1949 edition, p.~447, as Eberhardt and Glymour cite it) is
therefore very weak: although for every positive \(\varepsilon\) some
finite stage puts the straight rule within \(\varepsilon\) of the limit
forever after, the length of that stage is unknown, and at no point do
we know whether we are near the truth
\cite[p.~375]{EberhardtGlymour2011ReichenbachReading}. The precise form
of this boundary, including the fact that no finite prefix certifies the
limit, is \cref{thm:p15-reichenbach-straight-rule-boundary}.

Reichenbach supplements the argument with the network of inductions,
claiming that higher-order inductions make convergence faster than
inductions on single sequences. He gives no definition of speed of
convergence; his intuition is that cross-inductions bring in findings
from other domains, and his use of higher-order probability statements
even suggests uniform convergence. Eberhardt and Glymour reply that no proof is
given, that the claim is false without further assumptions, and that if
uniform convergence is not the intended sense it is unclear what faster
pointwise convergence buys. In a footnote they suggest this may be why
Reichenbach did not care when to use the rule of induction and when a
Bayesian update: both converge to the same limit, though one may worry
about what happens before convergence
\cite[p.~375 and n.~12]{EberhardtGlymour2011ReichenbachReading}.

If there is no limit, Reichenbach argues, no procedure could yield
inductive knowledge. He compares his method with a clairvoyant who
claims to know the limiting distribution (Eberhardt and Glymour cite the 1949
edition, p.~476): whether the person is clairvoyant can be checked only
by induction, and the straight rule at least guarantees pointwise
convergence while the self-declared clairvoyant may simply be wrong. Eberhardt and Glymour judge none of these arguments compelling, calling them
hand-waving at the crucial points \cite[p.~376]{EberhardtGlymour2011ReichenbachReading}.
The French original's version of the clairvoyant argument is reported in
\cref{sec:p15-reichenbach-straight-rule}.

Asked instead whether the procedure is unique in what it achieves, the
answer is no, and Reichenbach discusses this (Eberhardt and Glymour cite the 1949
edition, p.~447 and section~88). Many procedures converge pointwise when
a limit exists: the straight rule plus an arbitrary added function that
itself converges suitably; arbitrary guesses up to some point followed by
the straight rule; the straight rule plus a function tending to 0 as the
sequence grows. Reichenbach's response, Eberhardt and Glymour say, is confusing. He
grants that such altered estimates may even converge faster in some
circumstances and are equally legitimate, and yet claims that the
straight rule is, in an unspecified sense, unbiased and functionally
simplest. In Eberhardt and Glymour's view neither line survives careful analysis,
leaving an extremely weak justification, which Reichenbach regarded as
the best available without deceiving ourselves about the width of the
inductive gap \cite[p.~376 and n.~13]{EberhardtGlymour2011ReichenbachReading}.
That the rival rules listed have the same limit is part of
\cref{thm:p15-reichenbach-straight-rule-boundary}.

At several places (for example the 1949 edition, p.~481, as cited) the
justification is more pragmatic. Probability is a guide to action, and
since no empirical knowledge is certain, the theory is offered as a
workable way of guessing which theories will prove useful; Reichenbach
even claims that knowing whether a data sequence actually has a limit
would make no difference to scientific practice. Eberhardt and Glymour ask why, in
that case, he did not work with probabilities in finite sequences, and
note that in a reply to Russell he did just that, calling all real
probabilities finite frequencies and the talk of limits a fiction that
justifies our inductive procedures
\cite[p.~376]{EberhardtGlymour2011ReichenbachReading}.

Finally Eberhardt and Glymour compare the justification with practice. Pointwise
convergence is often not thought adequate: scientists frequently assume
much more, such as Gaussian or other parametric forms, which give
convergence rates and confidence intervals; Reichenbach might answer that
they deny the limits of their knowledge and that his logic explains why.
Recent non-parametric methods assume more than Reichenbach but still so
little that they support only asymptotic normality, so confidence
intervals concern only the limiting distribution, and without known
rates the question of what finite data license returns. Large
simulations now used to gauge those rates can be seen as supplying
``appraised posits'' about convergence; Eberhardt and Glymour close with the
remark that Reichenbach might feel vindicated
\cite[pp.~376--377]{EberhardtGlymour2011ReichenbachReading}. These are
general observations about statistical practice, offered without a
survey design or quantitative evidence.

\subsection{Critics: Popper, Nagel and Russell}
\label{sec:p15-w02b-p15-030-eg-critics}

Eberhardt and Glymour note that criticism came from every side, that Nagel's 1936
critical notice of the German \emph{Theory of Probability} in \emph{Mind} is
perhaps the most detailed and concise summary, that several points had
already been made by Popper, Tarski, Hertz and others, and that some
reappear in Russell's \emph{Human Knowledge} (1948). Because assigning
each point to one critic would mislead, they organize the material by
three foci: Popper on probabilities of theories, Nagel on the straight
rule, and Russell on the logical foundation
\cite[p.~377]{EberhardtGlymour2011ReichenbachReading}.

\paragraph{Popper.} In \emph{Logik der Forschung} (1934), before the
fullest exposition appeared, and in a 1935 notice in \emph{Erkenntnis},
Popper rejects the logic. He agrees that theories cannot be shown true,
but denies that they can even be given positive probability, and sees no
way for Reichenbach to avoid either an a priori foundation of
probability, which Reichenbach would reject, or an infinite regress of
higher-order probabilities, which Popper thinks cannot fix numerical
values \cite[p.~377]{EberhardtGlymour2011ReichenbachReading}.

Reichenbach offered two routes to the probability of a theory. The
first takes a sequence of experimentally testable singular consequences
of the theory and uses the frequency of confirmed ones. The second puts
the theory in a sequence of theories of the same reference class and uses
the frequency of true theories in it. Popper calls the second ludicrous:
there is no unique reference class of theories, even granted enough
theories; and if the theories in it can be judged true or false to
compute the frequency, a probability of the theory is redundant anyway.
Popper doubts that single theories can be judged true at all; saying
that Newton's law is true is not like saying that a coin landed heads.
If theories are judged only probabilistically, the regress returns, and
since each probability in an infinite regress is below 1, the
probability at the first level must be zero
\cite[pp.~377--378]{EberhardtGlymour2011ReichenbachReading}.

This point was pressed by many, including Russell, who took it to show
that two-valued logic is more fundamental than probability logic, and
C.~I. Lewis, who debated it with Reichenbach in the 1940s and 1950s with
respect to the foundations of epistemology. Reichenbach's reply in a
letter to Russell (published 1978), which Eberhardt and Glymour call opaque, seems
to be that higher-order claims are not probabilistically independent, so
their joint probability is not the product of their probabilities and
need not tend to 0. In the letter, \(a\) states that the event has
probability \(3/4\) and \(b\) states that \(a\) has probability \(1/2\);
to obtain the event's probability one also needs its probability if \(a\)
is false, which may exceed \(3/4\); the product computed by the critic is
the probability of the whole infinite conjunction of propositions on all
levels, which is indeed 0 (the letter refers to pp.~316--317 of the
German edition) \cite[p.~378]{EberhardtGlymour2011ReichenbachReading}.
By total probability, if the event has probability \(x\) when \(a\) is
false, its probability is \(\tfrac12\cdot\tfrac34+\tfrac12x=\tfrac38+\tfrac x2\),
which ranges over \([3/8,7/8]\) and exceeds \(3/4\) exactly when \(x>3/4\).
The general form of this computation, and the fact that an infinite
product of factors below 1 can be positive, are
\cref{thm:p15-reichenbach-higher-order-boundary}.

Popper and Nagel regard the first route as hopeless too, doubting that a
theory can be represented as an infinite conjunction of individually
testable singular statements. They and others add that on this route a
theory refuted by 10\% of its predictions would be true with probability
0.9, whereas it would more likely be counted false. Reichenbach's
response to Nagel's version (1939, and the 1949 edition, p.~436, as Eberhardt and Glymour cite them) does not, in Eberhardt and Glymour's judgment, meet the point.
He says his account gives a consistent interpretation of what scientists
might mean by the probability of a hypothesis, and in the English
edition explains: when the limits of exactness are drawn narrowly,
universal statements in science always have exceptions; when they are
drawn widely, a single exception is treated as incompatible with the
universal statement, for two reasons. Such statements usually have so
high a probability that one exception is a noticeable loss; and one
exception proves the statement false, and one dislikes using a
statement known to be false as a schematization. Eberhardt and Glymour find this
puzzling from someone who takes all empirical claims to be
probabilistic: if the probability of a theory is the frequency of its
positive instances, rather than the result of a Bayesian update, one
counterexample among many confirmations should change it only slightly.
For Popper, any third route, such as posits, is subjective
\cite[pp.~378--379]{EberhardtGlymour2011ReichenbachReading}.

\paragraph{Nagel and Hertz.} Nagel (1936) praises the precision of the
frequency interpretation but denies that it and the logic suffice to
represent and justify inductive inference, or describe it adequately.
Probability statements are unverifiable because they rest on sequences
whose limit, and even whose having a limit, is unknown. Nagel is not
persuaded by ``inductive verifiability''; he notes the lack of a
mathematical proof that higher-order probabilities or cross-inductions
speed convergence, and sees no value in a convergence guarantee when the
point of convergence is unknown. To a similar objection by Hertz (1936)
Reichenbach (1936) only repeats that his procedure is the best possible
and that nothing does better \cite[p.~379]{EberhardtGlymour2011ReichenbachReading}.

\paragraph{Creary's assessment.} A detailed later study of what can be
proved for Reichenbach's justification and his higher-order inductions
is Creary's 1969 Princeton thesis. Eberhardt and Glymour quote his four
conclusions (ch.~5, p.~129 there); in paraphrase, with Eberhardt and Glymour's
bracketed glosses:
\begin{enumerate}
\item The method of correction by cross-inductions (MC) gives no reason
  for choosing the lattice (for example, the grid of gravitational
  measurements from different planets) into which a given sequence is
  placed.
\item That MC is superior to Reichenbach's rule of induction (RRI\('\),
  the straight rule in a slightly modified form, p.~115 there) does not
  justify choosing MC over other simpler lattice-convergent rules.
\item MC's superiority, or even its parity, relative to RRI\('\) does not
  follow from Reichenbach's convergence theorems (1)--(4) (1949 edition,
  pp.~466--467; Creary p.~119).
\item Those theorems themselves rest on assumptions that would deprive
  any result proved with their help of the justificatory force
  Reichenbach intended.
\end{enumerate}
Eberhardt and Glymour call the assessment negative throughout and a vindication of
Nagel \cite[p.~380 and n.~14]{EberhardtGlymour2011ReichenbachReading}.

Nagel also denies that the logic describes scientific practice:
statements are judged probable not by a formal account but because no
plausible alternative exists, because the theory is aesthetically
attractive, or because it largely agrees with the evidence. He holds that
there are several more or less formal notions of probability and that
human factors enter into judging how probable a theory is
\cite[p.~380]{EberhardtGlymour2011ReichenbachReading}.

\paragraph{Tarski, Nagel and extensionality.} Nagel takes up an
objection of Tarski (1935): do the probabilities that replace truth
values express a syntactic relation among statements, or a semantic
relation between statements and what they describe? Nagel finds both
characterizations in probability implication. Eberhardt and Glymour suggest the
source of the dispute is that the value of a composite statement depends
on a third, conditional probability fixing the ``coupling'' of its
components; based on subsequence selection, it seems to add a syntactic
constraint to the semantic relation. Nagel and Tarski describe this as
the logic's not being ``extensional''. Reichenbach (1939) accepts that it
is not extensional in their strict sense (compositional, in Eberhardt and Glymour's
term) but denies that it is therefore intensional, that is, dependent on
the content of what statements describe; he presents the conditional
probability as the natural extension of the old dichotomy that
continuous-valued logics require. Eberhardt and Glymour doubt this would have
satisfied Nagel or Tarski, and find the standards in the debate and their
application to probabilistic logics unclear
\cite[pp.~380--381 and n.~15]{EberhardtGlymour2011ReichenbachReading}.

\paragraph{Russell.} Russell (1948) shares the worry about the logic's
status for another reason. Using the probability that an integer chosen
at random is prime (\emph{Human Knowledge}, pp.~366--368, as cited), he
shows that by rearranging the integers Reichenbach's definition can be
made to give any value between 0 and 1. Reichenbach, replying (1978),
accepts the conclusion and counts order-dependence among his theory's
advantages over class-based theories such as Keynes's. Russell takes the
dependence to show that the foundation cannot be stated in purely logical
terms, since order introduces a contingent, non-logical element
\cite[p.~381]{EberhardtGlymour2011ReichenbachReading}.

\begin{hbtheorem}{thm:p15-w02b-p15-030-russell-primes}{proved=annex:p15-w02b-p15-030-russell-primes}{Russell's rearrangement of the integers}
For every \(r\in[0,1]\) there is an enumeration \(m_1,m_2,\ldots\) of the
positive integers (each occurring exactly once) such that
\[
 \lim_{N\to\infty}\frac{|\{i\leq N: m_i\text{ is prime}\}|}{N}=r.
\]
For \(0<r<1\) one may take the enumeration that at step \(N+1\) lists the
least unused prime if fewer than \(rN\) primes have been listed so far,
and the least unused non-prime otherwise; then the count of primes among
the first \(N\) terms differs from \(rN\) by at most 1.
\end{hbtheorem}

The theorem uses only that there are infinitely many primes and
infinitely many non-primes; it therefore applies to any infinite,
co-infinite set of integers, which is the structural point behind
Russell's example.

Russell doubts that probability logic could be fundamental. In his words
as Eberhardt and Glymour quote them (\emph{Human Knowledge}, p.~368f):
\enquote{Statistical probability can only be estimated on a basis of
certainty, actual or postulated.} He adds that even in Reichenbach's
version two-valued truth plays the basic role, so the goal is not
reached \cite[p.~381]{EberhardtGlymour2011ReichenbachReading}.
Eberhardt and Glymour summarize Russell's view as three horns. If probability
values depend on the limit of an extensionally given sequence, they are
never attainable; if on a hierarchy of higher-order probabilities, these
cannot all be determined; if on posits, the posits depend on the order of
events. Russell did not share Reichenbach's indifference to the truth of
posits and probably did not accept his reply that blind posits are
justified as means to an end and require no belief in their truth (1978,
p.~407, as cited). For Russell, infinite regresses, hypothetical posits
and infinite extensional sequences cannot ground a logical calculus. He
proposes instead that the logic be built on finite sequences: scientific
probabilities concern large but finite populations anyway, finiteness
would allow precise estimation, and probability statements would then be
true or false, restoring two-valued logic at the foundation
\cite[pp.~381--382]{EberhardtGlymour2011ReichenbachReading}.

\paragraph{Nagel's translation dilemma.} Here Russell partly repeats
Nagel, who saw the rule of induction not as progress over the gap but as
a restatement of the problem. Since Reichenbach defines probability in
two-valued terms and places no limit on the complexity of statements,
the whole probability logic should be restatable in two-valued logic. If
the translation is possible, what is the probability logic for? If not,
in what sense does its treatment of induction solve the traditional
two-valued problem? Reichenbach (1935) distinguishes formal probability
logic, reducible to two-valued logic, from applied probability logic, for
which two-valued logic is at best an approximation, and compares the
situation to representing non-Euclidean space in Euclidean space; Eberhardt and Glymour say the analogy does not make things clearer to them
\cite[p.~382]{EberhardtGlymour2011ReichenbachReading}.

\subsection{Reichenbach on rival programs and standard problems}
\label{sec:p15-w02b-p15-030-eg-rivals}

\paragraph{Popper's asymmetry.} Reichenbach thought Popper in denial
about practice and about the consequences of his own theory (1935). One
may either respect the non-definitive character of scientific results,
which requires probabilities, or schematize, using only 0 and 1. In the
second case theories can be both verified and falsified; in the first,
neither can be done with certainty. Either way Popper's asymmetry between
verification and falsification disappears
\cite[p.~382]{EberhardtGlymour2011ReichenbachReading}.

He also regarded Popper's corroboration and severity of testing as
probability in disguise, too imprecise to be a metric, and compared
Popper to a fruit vendor who stacks the good fruit at the front of the
stall, fills the bags from the back, and denies that his method has
anything to do with the differing quality of the fruit. Popper sorts
unfalsified theories by a method he does not show to be reliable and
dismisses attempts at justification as ``metaphysical beliefs'', a charge
neither Popper nor Reichenbach wanted applied to his own work
\cite[p.~383]{EberhardtGlymour2011ReichenbachReading}.

\paragraph{Carnap, Hempel, Helmer and Oppenheim.} Against the
``deductive'' approaches of Carnap and Hempel (1949 edition,
pp.~456--461; see also Reichenbach 1935), Reichenbach argues that they do
not supply a space adequate for a probability metric; Eberhardt and Glymour call
this ironic, since Reichenbach did not specify a sigma-field either.
Carnap's probabilities rest on an a priori foundation, which leads
Reichenbach to say the result is hard to distinguish from a priori
methods such as the principle of indifference (p.~456). Eberhardt and Glymour are
unsure what he meant, since his usual objection to indifference concerns
its subjectivity and his usual objection to a priori methods concerns
their justification; either way the determination of the metric is at
issue. Against Hempel, Helmer and Oppenheim, who grant that no a priori
grounds select a metric but use initial data and a maximum-likelihood
assumption to fix it, Reichenbach objects that this assumes the data are
independent (p.~456); his network of inductions, integrating many
sequences in advanced knowledge, tests that assumption instead of making
it, and he counts the revisability of assumptions as an advantage (see
p.~464) \cite[p.~383]{EberhardtGlymour2011ReichenbachReading}.

\paragraph{Fisher and inverse probability.} Methods based entirely on
maximum likelihood answer the wrong question. In criticizing Fisher (1949
edition, p.~454f) Reichenbach says what matters is the inverse
probability of the hypothesis given the data, not its likelihood;
maximizing likelihood misleads because prior probabilities can reverse
the ordering, and once inverse probabilities are known likelihoods can be
set aside. He believed induction by enumeration would supply the priors
\cite[p.~383]{EberhardtGlymour2011ReichenbachReading}.

\paragraph{Ravens.} Reichenbach thought the inductive paradoxes vanish
once induction is treated probabilistically. He took Hempel's raven
paradox (1949 edition, p.~434f) to misapply contraposition to
probabilistic inference: ``if something is a raven, it is black'' may be
highly probable without ``if something is non-black, it is a non-raven''
being so; by Bayes's rule this depends on the base rates of ravens and of
black things. Eberhardt and Glymour describe this as a Bayesian view of
confirmation that dissolves the paradox for him
\cite[pp.~383--384]{EberhardtGlymour2011ReichenbachReading}.

\begin{hbtheorem}{thm:p15-w02b-p15-030-ravens}{proved=annex:p15-w02b-p15-030-ravens}{Probabilistic implication does not contrapose}
Let \(P\) be a probability with \(0<P(R)<1\) and \(0<P(\neg B)\), where
\(R\) is ``raven'' and \(B\) is ``black''. Then
\[
 P(\neg R\mid\neg B)=1-\frac{P(R)\,\bigl(1-P(B\mid R)\bigr)}{P(\neg B)}.
\]
In particular \(P(B\mid R)\) may be arbitrarily close to 1 while
\(P(\neg R\mid\neg B)=0\): with \(P(R)=0.9\), \(P(B\mid R)=0.9\) and
\(P(B\mid\neg R)=1\), one has \(P(\neg B)=0.09\) and
\(P(\neg R\mid\neg B)=0\). Conversely \(P(\neg R\mid\neg B)\) is close to
1 whenever \(P(R)(1-P(B\mid R))\) is small relative to \(P(\neg B)\).
\end{hbtheorem}

\paragraph{Grue.} Reichenbach's treatment of Goodman's riddle (1949
edition, p.~448f) is equally brief. ``Grue'' is a poor basis for
induction because it breaks the rule of using the narrowest reference
class. He grants that the rule of induction does not protect against
false posits in the short run: if ``grue'' means green until \(t\) and
blue afterwards, and all elements so far are green, ``All emeralds are
grue'' is confirmed, and generates false posits, up to \(t\) and is then
disconfirmed; in the long run one converges to the truth. But in
advanced knowledge, he argues, the inference is inferior because it
violates the rule \enquote{Use the narrowest common reference class
available}: by definition the property grue coincides with not-green
from the \((n+1)\)-st element on, and not-green, being the narrower
class, should be the basis of the inference
\cite[p.~384]{EberhardtGlymour2011ReichenbachReading}.

Eberhardt and Glymour's rebuttal is that up to the \(n\)-th item grue and green have
identical statistics, so no one could know in advance which to choose.
Contrary to what Reichenbach says elsewhere, ``narrowness'' here has
nothing to do with stable statistics and must rest on some other,
unstated criterion for individuating events. Goodman's riddle asks what
should constitute a reference class at all, that is, which elements
belong in a sequence used for a probability judgment; Reichenbach's
answer, in their view, misses this point
\cite[p.~384]{EberhardtGlymour2011ReichenbachReading}. Goodman's own
presentation is treated in
\cref{sec:p15-goodman-entrenchment-and-projectibility}.

\subsection{Eberhardt and Glymour's commentary}
\label{sec:p15-w02b-p15-030-eg-commentary}

Eberhardt and Glymour call the inductive logic a strange mixture of mathematical
precision and evasion of detail. Its aims are clear. Reichenbach wants a
probability logic that
\begin{enumerate}[label=(\alph*)]
\item is objective, hence the frequency interpretation;
\item is not a priori, hence the empirical focus;
\item gives a rational justification of inductive reasoning in science,
  hence the straight rule; and
\item is sensitive to the uncertainties in science, hence the many levels
  of probability.
\end{enumerate}
\cite[pp.~384--385]{EberhardtGlymour2011ReichenbachReading}

Their verdict on the machinery: beyond the frequency formalization of
probability there is no real need for all the logical apparatus, which,
if anything, makes the account heavier and more confusing. The main gap,
as others have said, is the justification by the straight rule; there
are no proofs that higher-order probabilities speed convergence, and it
is doubtful such proofs exist without substantial further assumptions
\cite[p.~385]{EberhardtGlymour2011ReichenbachReading}.

They also doubt that the logic captures scientific uncertainty, because
it represents three different kinds by one kind of probability:
uncertainty from noisy data (the subject of the theory of errors);
uncertainty because finitely many measurements underdetermine a lawlike
relation even without noise; and uncertainty in the physical quantity
itself, as on a metaphysical reading of Heisenberg's principle.
Reichenbach considers all three at different points. The first two are
epistemic and the third metaphysical: with direct access to the world
only the third would remain, and science tries to reduce the first two.
Traditional inductive logics addressed only the second; Reichenbach
thought them wishful because they treated data as certain, and he held
that no empirical knowledge is certain. But his logic does not separate
the contributions of the different sources
\cite[p.~385]{EberhardtGlymour2011ReichenbachReading}.

\begin{hbexample}{def:p15-w02b-p15-030-eg-noise-free}{Noise-free data on a line}
Suppose finitely many noise-free data points \((x,y)\), free of
metaphysical uncertainty, all lie on the line \(ax+b\). The true law
must pass through every point, but this does not determine it, because
anything may happen where there are no data. Which functional form is
best confirmed, most probable or best justified, and in what sense?
\cite[pp.~385--386]{EberhardtGlymour2011ReichenbachReading} That the data
leave a one-parameter family of distinct global fits for any finite
design is \cref{thm:p15-reichenbach-finite-fit}.
\end{hbexample}

\begin{hbdefinition}{def:p15-w02b-p15-030-eg-pf}{The \(k\)-dataset probability \(p_f\) of a function}
On Eberhardt and Glymour's account Reichenbach does not simply say that the data
imply, or make uniquely rational, the simplest fitting function. He holds
that each function \(f\) through the data points has an objective
probability \(p_f\) of being the true function, derived from \(k\) other
data sets: if \(f\) describes the data in \(k\cdot p_f\) of those data
sets, \(p_f\) is the probability that \(f\) is true. A footnote leaves
aside whether the data sets are independent samples
\cite[p.~386 and n.~16]{EberhardtGlymour2011ReichenbachReading}.
\end{hbdefinition}

So far there is no built-in preference for simpler theories. But
Reichenbach cannot, and seems not to think he can, avoid such a
preference altogether, because several functions fit all \(k\) data sets
yet differ where none of them has a point; here he does fall back on a
Bayesian account favouring simpler theories, without further argument.
Even if this handled underdetermination, noise and metaphysical
uncertainty must still be brought in. With suitable data, differential
likelihoods might separate hypotheses that constrain parameters
representing metaphysical uncertainty from hypotheses that do not; but
once noise is admitted, identifying and allocating uncertainty is in
principle massively underdetermined. Scientists hope for circumstances in
which the sources can be separated; Reichenbach's formalism makes no
attempt at this, although he knew Heisenberg's results and the theory of
errors well. The clean approach, Eberhardt and Glymour say, represents the sources
explicitly and separately, as Jeffrey conditionalization (1983) does: it
keeps noise in the data apart from uncertainty about functional form,
makes their confounding explicit, and shows how assumptions can separate
them. Reichenbach did not know it, but they guess he would have objected
to the many distributions it requires, which he would want to be
objective and for which an objective foundation seems far-fetched
\cite[p.~386]{EberhardtGlymour2011ReichenbachReading}.

Their closing assessment: their account will not have done justice to
every aspect; they doubt that a clean account of the probability of
theories can be given in Reichenbach's logic or in any purely formal
system; they have been critical of the justification of the straight
rule. Yet they value his starting point in actual scientific reasoning,
which offers a proper opposition to Popper's falsificationism and avoids
the quagmire of logical confirmation theory. Most of the logic is now of
historical interest, though some of his ideas about objective
probabilities in a Bayesian setting survive in philosophy, and ideas
similar to his mathematical theory have been developed, with proper
formalism, in modern statistics
\cite[p.~387]{EberhardtGlymour2011ReichenbachReading}.

The chapter ends with a note on its sources, which fixes the textual
basis of the whole reconstruction above. The account rests mainly on
\emph{The Theory of Probability}, in the German edition and the English
edition (the chapter's labels 1935c and 1949c), together with the papers
published before it in \emph{Erkenntnis} and elsewhere (labels 1931b,
1932b, 1934 and 1935e). The development of Reichenbach's foundations of
probability is traced through publications across his life, chiefly the
items labelled 1915, 1920, 1925, 1929, 1932a and 1933, with 1930 and 1931a
as further references. \emph{Experience and Prediction} (label 1938a) is
used to assemble the overall picture. His replies to critics in various
journals, chiefly the items labelled 1935d, 1936, 1939b, 1940 and 1949a,
and also 1938b, 1939a, 1941 and 1948, are used to settle points left
unclear in the systematic presentations. The authors close with an
explicit qualification: on some aspects of the view they
\enquote{remained at loss with regard to what exactly Reichenbach had in
mind} \cite[p.~387]{EberhardtGlymour2011ReichenbachReading}. The
reconstruction is therefore offered as their reading of a system whose
intended content they could not determine at every point, which matches
their earlier finding that the notations of the axioms and of the rule
of induction are never fully specified and that the joint syntax of the
two systems is never stated.

\begin{proofboundary}{def:p15-w02b-p15-030-eg-boundary}{What the Eberhardt--Glymour reconstruction establishes}
The theorems in this section are mathematical reconstructions with
stated hypotheses: the mode/mean contrast for a flat prior, Russell's rearrangement, and the failure of probabilistic
contraposition. None of them is a theorem of Reichenbach's own system,
whose syntax and semantics, as Eberhardt and Glymour stress, were never fully
specified; in particular nothing here proves the faster-convergence claim
for the network of inductions, Reichenbach's convergence theorems
(1)--(4) of the 1949 edition, or van Fraassen's general result, none of
which is reproduced in the chapter. The historical and philosophical
claims are reported with the chapter's page locators; quotations of
Reichenbach and Russell are given as the chapter quotes them, and
Creary's four conclusions are paraphrased from the chapter's quotation.
\end{proofboundary}

% ---------------------------------------------------------------------
\section{Niiniluoto on the development of the Hintikka program}
\label{sec:p15-w02b-p15-030-nn}

\subsection{The opening thesis and the research-programme framing}
\label{sec:p15-w02b-p15-030-nn-thesis}

Niiniluoto opens with the 1964 Jerusalem congress, where Hintikka
lectured ``Towards a Theory of Inductive Generalization''; two years
later Hintikka published a two-dimensional continuum of inductive
probability measures, and ten years later, with Niiniluoto, an axiomatic
system with \(K\geq2\) parameters. The thesis: these results showed once
and for all that there are systems of inductive logic in which genuine
universal generalizations have nonzero probability in an infinite
universe. Hintikka thereby disproved Popper's claim that inductive logic
is inconsistent and decisively improved Carnap's attempts to build
inductive logic as a theory of partial logical implication. The measures
later found applications in semantic information, confirmation and
acceptance, cognitive decision theory, analogy, truthlikeness and machine
learning, which, against Lakatos's early evaluation, confirm the
progressive character of the program
\cite[p.~311]{Niiniluoto2011HintikkaProgramReading}.

The framing is Lakatos's. In 1968 Lakatos proposed treating Carnap's
inductive logic as a methodological research programme, with a hard core
of assumptions, a positive heuristic for building a refutable protective
belt, and progress judged by its problem shifts. In a paper of 1969
(published 1974) he declared Popper's victory over a priori probabilistic
inductive logic complete, while calling inductive logic a degenerating
programme that was nevertheless a booming industry; David Miller (1994)
still holds a similar view. A footnote adds that Hintikka agrees with
Lakatos if the stress is on ``a priori'', but that Hintikka's
``victory'' over Carnap is entirely different from Popper's. Niiniluoto
argues the opposite of Lakatos's verdict: Hintikka's account of
generalization was a meeting point of several traditions and, as a
progressive turn, opened paths in Bayesian methodology, epistemology and
philosophy of science, with a potential in artificial intelligence that
is still largely unexplored
\cite[pp.~311--312 and n.~1]{Niiniluoto2011HintikkaProgramReading}.

\subsection{Prehistory and Carnap's development}
\label{sec:p15-w02b-p15-030-nn-prehistory}

Inductive logic goes back to the birth of the probability calculus in
the mid-seventeenth century, when probabilities were read as relative
frequencies, objective degrees of possibility and degrees of certainty
(Hacking 1975). For classical Bayesians such as the determinist Laplace,
probability was relative to ignorance of true causes; Jevons defended
probabilities as rational degrees of belief in the nineteenth century.
The Bayesian reading returned in the twentieth century in two forms: the
Cambridge school (W.~E. Johnson, Keynes, Broad, Jeffreys), which took
\(P(h/e)\) to be a logical relation between a hypothesis \(h\) and
evidence \(e\), with related logical views in Bolzano, Wittgenstein and
Waismann; and the subjective school (Ramsey, de Finetti), which read
degrees of belief as coherent betting ratios
\cite[p.~312]{Niiniluoto2011HintikkaProgramReading}.

In Finland, Eino Kaila published on ``probability logic'' in the spirit
of logical empiricism in the 1920s. His student G.~H. von Wright wrote a
1941 dissertation on the logical problem of induction and, influenced by
Keynes and Broad, tried to define the probability of a universal
statement through relative frequencies of properties. Von Wright was
Hintikka's most important teacher, and Hintikka and his students
continued this Finnish school. Carnap came late to the debate. Popper had
rejected induction in 1934 and criticized Reichenbach's frequentist
probability logic, with which Kaila sympathized; in a letter to Kaila of
28 January 1929 Carnap preferred a positivist solution in which
probability inferences are as analytic as syllogistic ones. Influenced
by Keynes, Jeffreys and Waismann, and to Popper's disappointment, Carnap
began his work on probability in 1942--44
\cite[p.~312]{Niiniluoto2011HintikkaProgramReading}.

The \emph{Logical Foundations of Probability} (LFP, 1950) developed the
measure \(c^*\), a generalization of Laplace's rule of succession; \emph{A
Continuum of Inductive Methods} (1952) introduced measures indexed by a
real parameter \(\lambda\), with \(c^*\) as one case. The case \(c^+\),
proposed by Bolzano and Wittgenstein, Carnap rejected because it allows
no learning from experience, a point Boole and Peirce had made against
Bayesian probabilities. With Kemeny, Carnap gave the
\(\lambda\)-continuum an axiomatic basis, and with Bar-Hillel he built a
theory of semantic information as part of his reply to Popper's praise of
improbability \cite[pp.~312--313]{Niiniluoto2011HintikkaProgramReading}.
Carnap's own 1950 measures are reconstructed in
\cref{sec:p15-carnap-logical-probability-and-induction}.

\subsection{The core assumptions and their reformulation}
\label{sec:p15-w02b-p15-030-nn-core}

\begin{hbdefinition}{def:p15-w02b-p15-030-nn-core}{Hacking's L1--L3, Carnap's C1--C7 and Hintikka's C3\('\), C6\('\)}
Following Hacking (1971), the typical assumptions of Carnap's inductive
logic are already in Leibniz:
\begin{description}
\item[(L1)] there is non-deductive evidence;
\item[(L2)] being a good reason for something is a relation between
  propositions;
\item[(L3)] there is an objective, formal measure of the degree to which
  one proposition is evidence for another.
\end{description}
They are also in Keynes's \emph{Treatise} (1921). Carnap's versions, in
the Preface to LFP, are:
\begin{description}
\item[(C1)] all inductive reasoning is reasoning in terms of probability;
\item[(C2)] inductive logic is probability logic;
\item[(C3)] inductive probability, or degree of confirmation, is a
  logical relation between two statements, a hypothesis and evidence;
\item[(C4)] the frequency concept is used in statistics but is unsuitable
  for inductive logic;
\item[(C5)] all principles and theorems of inductive logic are analytic;
\item[(C6)] the validity of induction does not depend on synthetic
  presuppositions.
\end{description}
LFP requires inductive logic to account for five kinds of inference:
\begin{enumerate}[label=(\alph*)]
\item direct inference, from a population to a sample;
\item predictive inference, from a sample to another sample;
\item inference by analogy, from one individual to another by known
  similarity;
\item inverse inference, from a sample to a population;
\item universal inference, from a sample to a universal hypothesis.
\end{enumerate}
Since Carnap did not claim \(c^*\) to be the only adequate explicatum
(LFP, p.~563), his method follows the heuristic principle
\begin{description}
\item[(C7)] use logic to distinguish the alternative states of affairs
  expressible in a formal language \(L\), then define inductive
  probabilities for sentences of \(L\) using symmetry assumptions about
  those states of affairs.
\end{description}
Hintikka replaces C3 and C6 by
\begin{description}
\item[(C3\('\))] inductive probability \(P(h/e)\) depends on the logical
  form of hypothesis \(h\) and evidence \(e\);
\item[(C6\('\))] inductive probabilities, and so inductive probabilistic
  inferences, may depend on extra-logical factors.
\end{description}
\cite[pp.~313--315]{Niiniluoto2011HintikkaProgramReading}
\end{hbdefinition}

Niiniluoto comments that the form \(P(h/e)\) ties Carnap to the Bayesian
tradition, that against the subjectivists Carnap wanted to remove all
``psychologism'' from inductive logic as Frege had from deductive logic,
that C4 accepts a dualism of physical and epistemic probability, and
that by C5 and C6 probability as partial entailment is independent of
all factual assumptions; in application, \(P(h/e)\) is computed relative
to the total evidence available. C1 and C2 leave open whether the basic
inductive notion is support (\(e\) confirms \(h\)) or acceptance (\(h\)
is rationally acceptable on \(e\)). Carnap held that Hume was right to
deny the validity of inductive inferences, so inductive logic evaluates
\(P(h/e)\), and the probabilities feed practical decisions through the
rule of maximizing expected utility (Stegmüller 1973). Normative decision
theory, Niiniluoto adds, contains ``quasi-psychological'' counterparts of
the ``purely logical'' inductive probabilities (Carnap 1971). Richard
Jeffrey followed Carnap in denying inductive acceptance rules (see the
debate in Lakatos 1968b). In the second
edition of LFP (1962) Carnap answered Popper by separating two senses of
``degree of confirmation'', the posterior probability \(P(h/e)\) and the
increase \(P(h/e)-P(h)\), a clarification of C3. C7 is what separates
the Carnapian program from Bayesianism in general, which admits any prior
\cite[pp.~313--314]{Niiniluoto2011HintikkaProgramReading}.

Hintikka relies on C7 in a new way that gives a new solution to problem
(e). His system satisfies C1, C2 and C4; but in reply to Mondadori (1987)
he insisted that his work amounted, for practical purposes, to a
refutation of Carnap's philosophical program, meaning the ``logicism'' of
C3, C5 and C6. C6\('\) allows contextual or ``local'' presuppositions
(Bogdan 1976), so that inductive probability involves pragmatic factors.
In what Hintikka calls ``logical pragmatics'', C3\('\) and C6\('\) are
combined with C7 so that the interplay of logical and extra-logical
factors is explicit. C5 then turns out to be ambiguous: some principles
of induction depend on pragmatic boundary conditions such as the
extra-logical parameters, while some mathematical theorems of inductive
logic are analytic \cite[pp.~314--315]{Niiniluoto2011HintikkaProgramReading}.
That distinction is the one kept in
\cref{def:p15-w02b-p15-030-claim-types}: the theorems below are
mathematical statements about measures with stated parameters; whether a
parameter value is appropriate is a separate, extra-logical question.

\subsection{Carnap's \(\lambda\)-continuum in Niiniluoto's presentation}
\label{sec:p15-w02b-p15-030-nn-carnap}

Probabilities in inductive logic are at least partly fixed by symmetry
assumptions on the language (Carnap 1962; Hintikka and Suppes 1966;
Niiniluoto and Tuomela 1973). In Carnap's \(\lambda\)-continuum a free
parameter \(\lambda\) gives the weight of logical, language-dependent
factors over empirical ones (observed frequencies) and so serves as an
index of caution in singular inference; Hintikka's 1966 system adds a
parameter \(\alpha\) governing how fast positive instances raise the
probability of a generalization \cite[p.~315]{Niiniluoto2011HintikkaProgramReading}.

\begin{hbdefinition}{def:p15-w02b-p15-030-nn-language}{Classification, state and structure descriptions, regular measures}
Let \(Q_1,\ldots,Q_K\) be a \(K\)-fold classification by mutually
exclusive predicates: every individual of the universe \(U\) satisfies
exactly one \(Q\)-predicate. Typically a finite monadic first-order
language \(L\) has basic predicates \(M_1,\ldots,M_k\) and each
\(Q\)-predicate is a conjunction \((\pm)M_1x\,\&\ldots\&\,(\pm)M_kx\), so
\(K=2^k\), and every predicate of \(L\) is a finite disjunction of
\(Q\)-predicates. Carnap generalized dichotomies \(\{M_j,\sim M_j\}\) to
families \(\mathbf M_j=\{M_{j1},\ldots,M_{jm}\}\) of exclusive
predicates, a \(Q\)-predicate choosing one member of each family (for
example a colour family and a discrete age family; Jeffrey 1980).

If \(L\) has individual names \(a_1,\ldots,a_N\) and is interpreted in a
universe of size \(N\) with each object named, a \emph{state
description} says, for each \(a_i\), which \(Q\)-predicate it satisfies;
a \emph{structure description} says how many individuals satisfy each
\(Q\)-predicate. Every sentence of \(L\) is a disjunction of state
descriptions; a structure description is the disjunction of the state
descriptions obtainable from one another by permuting individual
constants. The \emph{range} \(R(h)\) is the set of state descriptions
entailing \(h\). A \emph{regular} measure \(m\) gives each state
description \(s\) a nonzero value \(m(s)\), and \(m(h)\) is the sum of
\(m(s)\) for \(s\in R(h)\). The regular confirmation function is
\[
 (1)\qquad c(h/e)=\frac{m(h\&e)}{m(e)}.
\]
\cite[pp.~315--316]{Niiniluoto2011HintikkaProgramReading}
\end{hbdefinition}

\begin{hbdefinition}{def:p15-w02b-p15-030-nn-lambda}{The representative function and its special cases}
Let \(e_n\) describe a sample of \(n\) individuals by \(Q\)-predicates
and let \(n_i\geq0\) be the number observed in \(Q_i\)
(\(n_1+\cdots+n_K=n\)). Carnap's \(\lambda\)-continuum sets
\[
 (2)\qquad c(Q_i(a_{n+1})/e_n)=\frac{n_i+\lambda/K}{n+\lambda},
\]
the \emph{representative function} of the measure, a weighted average of
the observed frequency \(n_i/n\) and the relative width \(1/K\) of
\(Q_i\). Further:
\[
\begin{aligned}
 &(3)\quad [c(Q_i(a_{n+1})/e_n)-n_i/n]\to0\ \text{ as }n\to\infty,\\
 &(4)\quad c(Q_i(a_{n+2})/e_n\&Q_i(a_{n+1}))>c(Q_i(a_{n+1})/e_n),\\
 &(5)\quad c^*(Q_i(a_{n+1})/e_n)=\frac{n_i+1}{n+K}\qquad(\lambda=K),\\
 &(6)\quad \frac{n+1}{n+2}\qquad(n_i=n,\ K=2).
\end{aligned}
\]
(3) is \emph{Reichenbach's Axiom} (Carnap and Jeffrey 1971; Kuipers
1978b), (4) \emph{Positive Instantial
Relevance}, (5) Carnap's \(c^*\), and (6) Laplace's rule of succession.
If the \(Q\)-predicates have relative widths \(q_i\) with
\(q_1+\cdots+q_K=1\), (2) is replaced by
\[
 (2')\qquad \frac{n_i+q_i\lambda}{n+\lambda},
\]
(Carnap and Stegmüller 1959; Carnap 1980), and (2) is the case \(q_i=1/K\)
\cite[pp.~316--317]{Niiniluoto2011HintikkaProgramReading}.
\end{hbdefinition}

Niiniluoto reports that (3) is known to imply (4), that \(c^*\) divides
probability evenly among structure descriptions, that Laplace obtained
(6) by taking all compositions of an urn of white and black balls to be
equally probable, and, in a footnote, that statisticians know (2) and
(2\('\)) as symmetric and non-symmetric Dirichlet distributions (Festa
1993), which
Skyrms observed have been extended to a continuum of values
\cite[pp.~316--317 and n.~2]{Niiniluoto2011HintikkaProgramReading}.

As printed on p.~316, the special cases read ``\(\lambda=\infty\) gives
Reichenbach's Straight Rule'' and ``\(\lambda=4\) would give'' the
Tractatus range measure, which spreads probability evenly over state
descriptions and makes (2) equal to the evidence-independent \(1/K\).
By (2) itself the straight rule \(n_i/n\) is the value at \(\lambda=0\),
and the evidence-independent value \(1/K\) is the limit
\(\lambda\to\infty\) (\cref{thm:p15-w02b-p15-030-lambda}); the remark on
p.~323 that A2 excludes \(\lambda=0\) agrees with this reading. The
printed pair is therefore read here as \(\lambda=0\) and
\(\lambda=\infty\).

\begin{hbtheorem}{thm:p15-w02b-p15-030-lambda}{proved=annex:p15-w02b-p15-030-lambda}{Elementary properties of the representative function}
Let \(K\geq2\), \(0<\lambda<\infty\) and \(0\leq n_i\leq n\).
\begin{enumerate}
\item For \(n\geq1\),
  \(\dfrac{n_i+\lambda/K}{n+\lambda}
   =\dfrac{n}{n+\lambda}\cdot\dfrac{n_i}{n}+\dfrac{\lambda}{n+\lambda}\cdot\dfrac1K\),
  and
  \(\Bigl|\dfrac{n_i+\lambda/K}{n+\lambda}-\dfrac{n_i}{n}\Bigr|\leq\dfrac{\lambda}{n+\lambda}\).
  Hence (3) holds along every sequence of samples with \(n\to\infty\).
\item (4) holds for the representative function: for all \(n\geq0\),
  \(\dfrac{n_i+1+\lambda/K}{n+1+\lambda}>\dfrac{n_i+\lambda/K}{n+\lambda}\).
\item As \(\lambda\to0\) the value tends to \(n_i/n\) (for \(n\geq1\)); as
  \(\lambda\to\infty\) it tends to \(1/K\). At \(\lambda=K\) it is (5),
  and (5) with \(n_i=n\), \(K=2\) is (6).
\end{enumerate}
\end{hbtheorem}

When the universe may grow without limit, \(c(h/e)\) is defined as the
limit of the finite-universe values (1) as \(N\to\infty\), and then every
measure of the \(\lambda\)-continuum gives probability zero to universal
generalizations on singular evidence. Carnap admitted this ``may seem
astonishing'', concluded that the role of universal sentences in science
had been overestimated, and proposed instead the \emph{instance
confirmation} of a law: the probability that a new individual not
mentioned in \(e\) satisfies it. (The 1950 counterpart of this contrast is
\cref{thm:p15-carnap-universal}.) In Lakatos's terms this reduction of
universal to singular predictive inference was a regressive problem
shift. Von Wright objected (1951). Following Keynes, he knew that a
universal generalization \(h\) is confirmed by positive singular
evidence \(e_n=i_1\&\ldots\&i_n\) entailed by \(h\) provided (i) its prior
is not minimal and (ii) new confirmations are not maximally probable
given the earlier ones \cite[p.~317]{Niiniluoto2011HintikkaProgramReading}.

\begin{hbtheorem}{thm:p15-w02b-p15-030-principal}{proved=annex:p15-w02b-p15-030-principal}{Von Wright's Principal Theorem of Confirmation (7)}
Let \(P\) be a probability, let \(h\) entail \(e_n\& i_{n+1}\), and
suppose \(P(e_n\&i_{n+1})>0\). If \(P(h)>0\) and \(P(i_{n+1}/e_n)<1\),
then
\[
 (7)\qquad P(h/e_n\&i_{n+1})=\frac{P(h/e_n)}{P(i_{n+1}/e_n)}>P(h/e_n).
\]
\end{hbtheorem}

Since Carnap's system fails (7) for laws, von Wright said it is no
confirmation theory at all (in the 1957 book, pp.~119 and~215, as
cited); nor does
it settle the Keynes--Nicod dispute over when the posterior of a law
tends to its maximum,
\[
 (8)\qquad P(h/e_n)\to1\quad\text{as }n\to\infty.
\]
Popper too argued for zero logical probability of universal laws.
Lakatos (1968) called the assumption \(P(h)>0\) for genuine universal
statements the ``Jeffreys--Keynes postulate'' and Carnap's dispensability
thesis the ``weak atheoretical thesis''. If \(P(h)=0\) for laws, then
\(P(h/e)=0\) for every \(e\), and the difference measure and Popper's
corroboration are trivialized too (Niiniluoto and Tuomela 1973,
pp.~212--216, 242--243). Instance confirmation confines
inductive logic to singular sentences; Mary Hesse (1974) proposed keeping
Carnap's framework for generalizations in finite universes, but this
makes inductive logic depend on a synthetic assumption about the size of
the universe, against C6, and gives qualitatively wrong results, the
strongest confirmation going to generalizations that allow many
\(Q\)-predicates even when the evidence concentrates on a few. In the
Schilpp volume (p.~977) Carnap said he had constructed confirmation
functions giving universal generalizations nonzero probability but that
they were ``considerably more complicated'' than the \(\lambda\)-system;
he never published them. Against this background Hintikka's Jerusalem
system of 1964 was striking \cite[pp.~317--318]{Niiniluoto2011HintikkaProgramReading}.

\subsection{Constituents and Hintikka's systems}
\label{sec:p15-w02b-p15-030-nn-constituents}

Hintikka's solution divides probability among constituents. He had met
the tool in von Wright's Helsinki lectures of 1947--48, where logical
truth in monadic predicate logic was characterized by distributive
normal forms: a tautology allows all constituents, mutually exclusive
descriptions of the constitution of the universe. Hintikka's insight of
1948, at twenty-one, was how to extend distributive normal forms to full
first-order logic with relations; this became his 1953 dissertation
\cite[p.~318]{Niiniluoto2011HintikkaProgramReading}.

\begin{hbdefinition}{def:p15-w02b-p15-030-nn-constituent}{Constituents, width and distributive normal form}
In a monadic language with \(Q\)-predicates \(Q_1,\ldots,Q_K\), a
\emph{constituent} \(C^w\) says which \(Q\)-predicates are nonempty and
which empty in \(U\):
\[
 (9)\qquad (\pm)(\exists x)Q_1(x)\,\&\ldots\&\,(\pm)(\exists x)Q_K(x).
\]
If \(Q_i\), \(i\in CT\), are exactly those claimed nonempty, (9) is
\[
 (10)\qquad \bigwedge_{i\in CT}(\exists x)Q_i(x)\ \&\ (x)\Bigl[\bigvee_{i\in CT}Q_i(x)\Bigr].
\]
The cardinality \(w\) of \(CT\) is the \emph{width}. \(C^K\), which says
that every \(Q\)-predicate is exemplified, is the \emph{atomistic}
constituent: if it is true, no universal generalization of \(L\) is
true, and \(U\) is atomistic with respect to \(L\). There are \(2^K\)
constituents, including the empty one of width zero, which is a
contradiction; the others are maximally consistent and complete theories
of \(L\), each specifying a ``possible world'' by primitive monadic
predicates, connectives and quantifiers. Constituents are mutually
exclusive and their disjunction is a tautology; with finitely many
individual constants each is a disjunction of state descriptions, or of
structure descriptions. Every consistent generalization \(h\) of \(L\) (a
quantified sentence without individual constants) has a
\emph{distributive normal form}
\[
 (11)\qquad \vdash h\equiv\bigvee_{i\in I_h}C_i,
\]
where \(I_h\) indexes the constituents that entail \(h\). Constituents are
the \emph{strong} generalizations of \(L\); the others are \emph{weak}
\cite[pp.~318--319]{Niiniluoto2011HintikkaProgramReading}.
\end{hbdefinition}

The source prints (11) with ``\(=\)''; the turnstile and the role of the
sign show that logical equivalence is meant, and it is written
\(\equiv\) here. By (11) the probability of any generalization reduces
to probabilities of constituents.

Let \(e\) describe a finite sample of \(n\) individuals in which \(c\)
different kinds (\(Q\)-predicates) are observed, written \(e^c_n\) when
needed. A constituent \(C^w\) is compatible with \(e^c_n\) only if
\(c\leq w\leq K\), and the compatible ones are \(C^c\) together with the
\(\binom{K-c}{i}\) constituents of width \(c+i\) for each
\(i=1,\ldots,K-c\). By Bayes's formula, writing \(C^{c+i}\) for any
compatible constituent of width \(c+i\) (all of which have equal prior
and likelihood in the systems below),
\[
 (12)\qquad P(C^w/e)=\frac{P(C^w)P(e/C^w)}
   {\sum_{i=0}^{K-c}\binom{K-c}{i}P(C^{c+i})P(e/C^{c+i})}.
\]
Priors \(P(C^w)\) and likelihoods \(P(e/C^w)\) must therefore be
specified \cite[p.~319]{Niiniluoto2011HintikkaProgramReading}.

\paragraph{Jerusalem and combined systems.} The Jerusalem system
(Hintikka 1965) divides probability evenly among constituents and then
evenly among the state descriptions belonging to each. The combined
system (also 1965) divides it evenly among constituents, then evenly
among the structure descriptions satisfying each, then evenly among the
state descriptions in each structure description. In both, every
constituent has prior \(1/2^K\), and exactly one constituent has
asymptotic probability one: the minimal constituent \(C^c\), which says
that \(U\) instantiates exactly the \(c\) kinds seen in the sample:
\[
 (13)\qquad P(C^c/e^c_n)\to1,\qquad P(C^w/e^c_n)\to0\ (w>c),
\]
as \(n\to\infty\) with \(c\) fixed.
Hence a constituent claiming that some uninstantiated kind exists tends
to probability zero, and for a weak generalization
\[
 (14)\qquad P(h/e^c_n)\to1\iff C^c\vdash h\qquad(n\to\infty,\ c\text{ fixed}),
\]
which answers the Keynes--Nicod question
\cite[pp.~319--320]{Niiniluoto2011HintikkaProgramReading}.

\begin{hbwarning}{def:p15-w02b-p15-030-nn-empty}{The empty constituent and the prior \(1/2^K\)}
Read literally, ``every constituent has prior \(1/2^K\)'' gives the empty
constituent, a contradiction, probability \(1/2^K\), which no probability
measure allows. On a strict reading the equal priors are \(1/(2^K-1)\) on
the consistent constituents. Every posterior formula below conditions on
evidence with \(c\geq1\), cancels the common normalizing constant of the
prior, and is unaffected by this choice.
\end{hbwarning}

\paragraph{The two-dimensional continuum.} In 1966 Hintikka found a
system containing his earlier measures and the whole
\(\lambda\)-continuum as special cases. Likelihoods given \(C^w\) are
computed as in Carnap's continuum but with the universe restricted to the
\(w\) kinds that \(C^w\) instantiates; for evidence compatible with
\(C^w\),
\[
 (15)\qquad P(Q_i(a_{n+1})/e\&C^w)=\frac{n_i+\lambda/w}{n+\lambda},
\]
\[
 (16)\qquad P(e/C^w)=\frac{\Gamma(\lambda)}{\Gamma(n+\lambda)}
   \prod_{j=1}^{c}\frac{\Gamma(n_j+\lambda/w)}{\Gamma(\lambda/w)},
\]
where \(\Gamma\) is the gamma function (\(\Gamma(n+1)=n!\)) and the
product runs over the \(c\) observed kinds. The prior of \(C^w\) is taken
proportional to the Carnapian probability that a set of \(\alpha\)
individuals is compatible with \(C^w\):
\[
 (17)\qquad P(C^w)=\frac{\Gamma(\alpha+w\lambda/K)/\Gamma(w\lambda/K)}
   {\sum_{i=0}^{K}\binom Ki\,\Gamma(\alpha+i\lambda/K)/\Gamma(i\lambda/K)}.
\]
If \(\alpha=0\), (17) gives equal priors,
\[
 (18)\qquad P(C^w)=1/2^K\ \text{ for all }C^w,
\]
and letting \(\lambda\to\infty\) then yields the Jerusalem system. A
small \(\alpha\) expresses strong a priori weight in generalization, as a
small \(\lambda\) does in singular inference. For \(0<\alpha<\infty\),
\[
 (19)\qquad P(C^w)<P(C^{w'})\iff w<w'.
\]
\cite[p.~320]{Niiniluoto2011HintikkaProgramReading}

The minimal constituent \(C^c\) says that the universe resembles the
sample; it is the simplest unrefuted constituent in the sense of
ontological parsimony. By (19) it is initially the least probable, yet by
(13) the only one whose probability tends to one on growing but similar
evidence \cite[p.~321]{Niiniluoto2011HintikkaProgramReading}.

\paragraph{The generalized combined system.} Choosing \(\lambda\) as a
function of width, \(\lambda(w)=w\), gives Hintikka's \emph{generalized
combined system}; the 1965 combined system is its case \(\alpha=0\). The
formulas become
\[
\begin{aligned}
 (15')\quad& P(Q_i(a_{n+1})/e\&C^w)=\frac{n_i+1}{n+w},\\
 (16')\quad& P(e/C^w)=\frac{(w-1)!}{(n+w-1)!}\prod_{j=1}^c(n_j!),\\
 (17')\quad& P(C^w)=\frac{(\alpha+w-1)!/(w-1)!}
   {\sum_{i=0}^K\binom Ki(\alpha+i-1)!/(i-1)!},\\
 (20)\quad& P(C^w/e^c_n)=\frac{(\alpha+w-1)!/(n+w-1)!}
   {\sum_{i=0}^{K-c}\binom{K-c}{i}(\alpha+c+i-1)!/(n+c+i-1)!}.
\end{aligned}
\]
When \(n=\alpha\), (20) reduces to \(1/\sum_{i=0}^{K-c}\binom{K-c}i=1/2^{K-c}\).
For \(n\) and \(\alpha\) large relative to \(K\), the approximation
\((m+n)!\simeq m!\,m^n\), valid for \(m\) large relative to \(n^2\) (Carnap
1950, p.~150), gives
\[
 (21)\qquad P(C^w/e^c_n)\simeq\frac{(\alpha/n)^{w-c}}{(1+\alpha/n)^{K-c}}
\]
(Niiniluoto 1987, p.~88), which displays the asymptotics (13). The
representative function is
\[
 (22)\quad P(Q_i(a_{n+1})/e^c_n)=(n_i+1)\,
 \frac{\sum_{i=0}^{K-c}\binom{K-c}{i}\frac{(\alpha+c+i-1)!}{(n+c+i)!}}
      {\sum_{i=0}^{K-c}\binom{K-c}{i}\frac{(\alpha+c+i-1)!}{(n+c+i-1)!}}.
\]
If \(h\) is a universal generalization saying that \(b\) given kinds are
empty, and \(h\) is compatible with \(e\),
\[
 (23)\quad P(h/e^c_n)=
 \frac{\sum_{i=0}^{K-b-c}\binom{K-b-c}{i}\frac{(\alpha+c+i-1)!}{(n+c+i-1)!}}
      {\sum_{i=0}^{K-c}\binom{K-c}{i}\frac{(\alpha+c+i-1)!}{(n+c+i-1)!}},
\]
\[
 (24)\quad P(h/e)\simeq\frac1{(1+\alpha/n)^b}
\]
for large \(\alpha\) and \(n\); (24) tends to one as \(n\) grows, in
agreement with (14). Conversely, as \(\alpha\to\infty\), (20) gives
\(P(C^w/e)\to1\) if and only if \(w=K\), and the same holds for priors:
\[
 (25)\qquad \alpha\to\infty\implies P(C^K)\to1\text{ and }P(C^w)\to0\ (w<K).
\]
\cite[pp.~321--322]{Niiniluoto2011HintikkaProgramReading}

In (22) the inner summation index is also printed \(i\), clashing with
the cell index; the cell index belongs only to \(n_i\). The formula
applies to kinds \(Q_i\) already observed in \(e\)
(\cref{thm:p15-w02b-p15-030-combined}).

Niiniluoto adds that, more generally, the probabilities of the
\(\lambda\)--\(\alpha\) continuum approach Carnap's as \(\alpha\to\infty\);
that (25) explains why all Carnap's measures give universal laws
probability zero, since they fix a priori, sceptically, that the
atomistic constituent has prior one; and that Carnap's continuum is the
only special case (\(\alpha=\infty\)) of the two-dimensional continuum in
which (13) fails \cite[p.~322]{Niiniluoto2011HintikkaProgramReading}.

\begin{hbtheorem}{thm:p15-w02b-p15-030-likelihood}{proved=annex:p15-w02b-p15-030-likelihood}{The likelihoods (16) and (16\('\)) from the predictive rule}
Fix \(w\geq1\) and \(0<\lambda<\infty\). Suppose that, given \(C^w\), the
probability that \(a_{m+1}\) falls in \(Q_i\) after a sample \(e_m\) of
\(a_1,\ldots,a_m\) compatible with \(C^w\) is
\((m_i+\lambda/w)/(m+\lambda)\) for each of the \(w\) kinds \(Q_i\)
claimed nonempty by \(C^w\), where \(m_i\) counts \(Q_i\) in \(e_m\). Then for
every sample description \(e\) of \(a_1,\ldots,a_n\) compatible with
\(C^w\), with counts \(n_j\) in the \(c\) observed kinds, \(P(e/C^w)\) is
given by (16). With \(\lambda=w\) it is (16\('\)).
\end{hbtheorem}

\begin{hbtheorem}{thm:p15-w02b-p15-030-combined}{proved=annex:p15-w02b-p15-030-combined}{Posterior behaviour of the generalized combined system}
Take the predictive rule (15\('\)) given each constituent, hence by
\cref{thm:p15-w02b-p15-030-likelihood} the likelihoods (16\('\)), and priors proportional to
\(\Gamma(\alpha+w)/\Gamma(w)\) for the constituents of width \(w\geq1\)
(this is (17\('\)), with factorials read as gamma values), where
\(0\leq\alpha<\infty\). Let \(e=e^c_n\) be a sample description with
\(n\geq c\geq1\). Then:
\begin{enumerate}
\item the posterior of each compatible constituent of width \(w\) is (20),
  with \((x-1)!\) read as \(\Gamma(x)\); when \(n=\alpha\) it equals
  \(1/2^{K-c}\);
\item (13) holds, and so does (14) for every generalization \(h\) with
  distributive normal form (11);
\item (23) holds for a universal generalization \(h\) compatible with
  \(e\) that declares \(b\) kinds empty;
\item (22) is the predictive probability of every kind \(Q_i\) observed
  in \(e\) (\(n_i\geq1\)).
\end{enumerate}
\end{hbtheorem}

\begin{hbtheorem}{thm:p15-w02b-p15-030-priors}{proved=annex:p15-w02b-p15-030-priors}{The prior assignments (17), (18), (19) and (25)}
\begin{enumerate}
\item In (17), \(\alpha=0\) makes the weight
  \(\Gamma(\alpha+w\lambda/K)/\Gamma(w\lambda/K)\) equal to \(1\) for every
  \(w\geq1\), so all consistent constituents receive equal priors.
\item For the priors (17\('\)) on constituents of width \(w\geq1\) and
  \(0<\alpha<\infty\): \(P(C^w)<P(C^{w'})\iff w<w'\), which is (19).
\item For the same priors, as \(\alpha\to\infty\), \(P(C^K)\to1\) and
  \(P(C^w)\to0\) for every \(w<K\), which is (25) for the generalized
  combined system.
\end{enumerate}
\end{hbtheorem}

The chapter states (19) and (25) for the general prior (17) with
arbitrary \(\lambda\), and states without proof that the whole
\(\lambda\)--\(\alpha\) continuum approaches Carnap's as
\(\alpha\to\infty\). The proofs here cover the generalized combined
system, where the gamma ratios differ by integer shifts; the general
case needs the asymptotics of \(\Gamma(\alpha+x)/\Gamma(\alpha+y)\) for
non-integer \(x-y\), which the chapter does not supply and which is not
reproduced here.

\begin{hbexample}{def:p15-w02b-p15-030-nn-worked}{A worked case of the generalized combined system}
Let \(k=2\) primitive predicates, so \(K=4\), and let a sample show
\(c=2\) kinds, \(Q_1\) and \(Q_2\), with \(n_1=12\), \(n_2=8\), \(n=20\).
Take \(\alpha=10\). By (20) the compatible constituents are \(C^2\)
(only \(Q_1,Q_2\) exist), two of width 3 and one of width 4, with
\[
 P(C^2/e)=\tfrac{253}{607}\approx0.417,\quad
 P(C^3/e)=\tfrac{138}{607}\approx0.227,\quad
 P(C^4/e)=\tfrac{78}{607}\approx0.129,
\]
the middle value for each of the two constituents of width 3.
Thus the universal generalization ``no individual is \(Q_3\) or
\(Q_4\)'' has probability \(253/607\), and ``no individual is \(Q_3\)''
has \((253+138)/607=391/607\approx0.644\); both agree with (23). By (22),
the probability that the next individual is \(Q_1\) is
\(1391/2428\approx0.573\), against the observed frequency \(0.6\). At
\(n=\alpha=10\) every compatible constituent has \(1/4\). With the same
proportions at \(n=100\), \(P(C^2/e)=1751/2189\approx0.800\), each width-3
constituent has \(206/2189\approx0.094\) and \(C^4\) has
\(26/2189\approx0.012\), illustrating (13); the approximation (21) gives
\(1/(1.1)^2\approx0.826\) for \(C^2\).
\end{hbexample}

\subsection{Axiomatic inductive logic}
\label{sec:p15-w02b-p15-030-nn-axioms}

Axiomatic inductive logic seeks general rationality principles that
narrow the admissible measures. W.~E. Johnson (1932) gave the first such
treatment; his main results were rediscovered independently by Kemeny and
Carnap in 1952--54 \cite[p.~322]{Niiniluoto2011HintikkaProgramReading}.

\begin{hbdefinition}{def:p15-w02b-p15-030-nn-a1a6}{Axioms A1--A6}
Let \(P\) be real-valued on pairs \((h,e)\) of sentences of a finite
monadic language \(L\), \(e\) consistent.
\begin{description}
\item[(A1)] \emph{Probability axioms.}
\item[(A2)] \emph{Finite regularity}: for singular \(h\) and \(e\),
  \(P(h/e)=1\) only if \(\vdash e\supset h\).
\item[(A3)] \emph{Symmetry with respect to individuals}: \(P(h/e)\) is
  invariant under every permutation of individual constants.
\item[(A4)] \emph{Symmetry with respect to predicates}: \(P(h/e)\) is
  invariant under every permutation of the \(Q\)-predicates.
\item[(A5)] \emph{\(\lambda\)-principle}: there is a function \(f\) with
  \(P(Q_i(a_{n+1})/e)=f(n_i,n)\).
\item[(A6)] \emph{\(c\)-principle}: there is a function \(f\) with
  \(P(Q_i(a_{n+1})/e^c_n)=f(n_i,n,c)\).
\end{description}
\cite[pp.~322--324]{Niiniluoto2011HintikkaProgramReading}
\end{hbdefinition}

For personalists A1 is the only general constraint; it makes
probabilities coherent betting ratios. A2 forbids prior probability one
for any contingent singular sentence. A3 is de Finetti's exchangeability (Carnap and Jeffrey 1971; Hintikka
1971):
the predictive probability depends on \(e\) only through
\(n_1,\ldots,n_K\), not on the order of observation. A4 makes the
\(Q\)-predicates symmetric, so \(P(Q_i(a_j))=1/K\). A5 is Johnson's
``sufficientness postulate'' (Carnap's ``axiom of predictive
irrelevance''): the representative function does not depend on the
counts \(n_j\), \(j\neq i\), in other cells, given their sum
\cite[pp.~322--323]{Niiniluoto2011HintikkaProgramReading}.

\begin{hbtheorem}{thm:p15-w02b-p15-030-kemeny-carnap}{proved=annex:p15-w02b-p15-030-kemeny-carnap}{The Kemeny--Carnap theorem on singular sentences}
Let \(P\) be defined on singular sentences about \(a_1,a_2,\ldots\) and
satisfy A1--A5, and suppose \(K\geq3\), or \(K=2\) and \(f(x,n)\) is
affine in \(x\) for each \(n\). Then either \(Kf(0,1)<1\) and, with
\[
 \lambda=\frac{Kf(0,1)}{1-Kf(0,1)}\in(0,\infty),\qquad
 f(n_i,n)=\frac{n_i+\lambda/K}{n+\lambda}\quad(0\leq n_i\leq n),
\]
or \(Kf(0,1)=1\) and \(f\equiv1/K\) (the limiting member
\(\lambda=\infty\)). Conversely, every member of this family,
\(0<\lambda\leq\infty\), is the predictive rule of a probability on
singular sentences satisfying A1--A5. The value \(\lambda=0\) is excluded
by A2.
\end{hbtheorem}

The theorem is Niiniluoto's statement made precise in three respects: it
concerns \(P\) on singular sentences (the extension to quantified
sentences in an infinite universe is Carnap's limit definition, described
above); it records the limiting member \(\lambda=\infty\), which satisfies
A1--A5; and it spells out the \(K=2\) proviso, which the chapter states as
requiring that \(f\) be a linear function of \(n_i\)
\cite[p.~323]{Niiniluoto2011HintikkaProgramReading}. The coherence of the
predictive family used in the converse is
\cref{thm:p15-carnap-history-family}; the history of the sufficientness
postulate is in \cref{sec:p15-carnap-logical-probability-and-induction}.

Dropping A4, Niiniluoto reports, gives the form (2\('\)); hence a regular,
exchangeable inductive measure is Carnapian exactly when it satisfies A5,
and Laplace's \(c^*\) satisfies A5. But A5 is very strong: it makes the
predictive probability independent of the variety of evidence, the number
\(c\) of cells with \(n_i>0\). Since the number of universal
generalizations falsified by \(e\) is also a simple function of \(c\), A5
makes induction purely enumerative and excludes its eliminative side
(Hintikka 1968). The representative function (22) depends on \(c\), and
the failure of the \(\lambda\)-continuum on generalization is an unhappy
consequence of A5 \cite[p.~323]{Niiniluoto2011HintikkaProgramReading}.

\paragraph{The \(K\)-dimensional system.} In 1974 (published 1976)
Hintikka and Niiniluoto replaced A5 by A6. The number \(c\) expresses the
variety of \(e\) and indicates how many generalizations it has falsified.
As Niiniluoto reports, they proved that measures satisfying A1--A4 and A6
form a \emph{\(K\)-dimensional system} determined by \(K\) parameters,
printed as
\[
 \lambda=\frac{Kf(1,K+1,K)}{1-Kf(1,K+1,K)}-1,\qquad
 \gamma_c=f(0,c,c)\quad(c=1,\ldots,K-1),
\]
with \(\lambda>-K\) and
\[
 (26)\qquad 0<\gamma_c\leq\frac{\lambda/K}{c+\lambda}.
\]
Kuipers (1978) calls this the NH-system. The upper bound in (26) is the
value \(\delta_c\) of \(f(0,c,c)\) in Carnap's continuum. For infinite
universes the atomistic constituent has
\[
 P(C^K)=\frac{\gamma_1\cdots\gamma_{K-1}}{\delta_1\cdots\delta_{K-1}},
\]
so \(P(C^K)=1\) iff \(\gamma_i=\delta_i\) for all \(i=1,\ldots,K-1\): Carnap's
continuum is the only special case of the \(K\)-dimensional system that
gives no universal generalization nonzero probability, and is biased in
assigning the atomistic constituent prior one
\cite[p.~324]{Niiniluoto2011HintikkaProgramReading}. (The product formula
is printed with a typographical slip, ``\(\gamma_1\ldots\gamma)K-1\)'';
the reading \(\gamma_1\cdots\gamma_{K-1}\) is fixed by the following
sentence.)

These results are reported from the chapter; their proofs are in
Hintikka and Niiniluoto (1976), which the chapter cites and does not
reproduce. What follows from the stated formulas is proved here.

\begin{hbtheorem}{thm:p15-w02b-p15-030-kdim}{proved=annex:p15-w02b-p15-030-kdim}{Consequences of the stated \(K\)-dimensional formulas}
\begin{enumerate}
\item In Carnap's continuum with \(0<\lambda<\infty\),
  \(\delta_c=f(0,c,c)=(\lambda/K)/(c+\lambda)\), and \(\delta_c<1/K\).
\item Given the product formula for \(P(C^K)\) and the bounds (26),
  \(P(C^K)=1\) iff \(\gamma_i=\delta_i\) for every \(i=1,\ldots,K-1\);
  otherwise \(0<P(C^K)<1\), so some constituent of smaller width has
  positive prior.
\item If the predictive rule (27) below holds for the constituent \(C^K\)
  (the only one compatible with evidence in which all \(K\) kinds occur),
  then \(f(1,K+1,K)=(1+\lambda/K)/(K+1+\lambda)\), and
  \[
   \lambda=\frac{Kf(1,K+1,K)}{1-Kf(1,K+1,K)}-K .
  \]
\end{enumerate}
\end{hbtheorem}

Part~3 does not match the printed parameter formula, whose last term is
\(-1\) rather than \(-K\); the two agree only when \(K=1\). Since the
1976 paper is not available here, the discrepancy is recorded rather than
resolved: either the printed formula or the parametrization implicit in
(27) differs from the original by a slip.

Niiniluoto continues: reducing all inductive probabilities to \(K\)
parameters about very simple singular predictions answers Stegmüller's
(1973) claim that betting on universal generalizations makes no sense,
since in this system a bet on a law is equivalent to a system of \(K\)
bets on singular sentences on finite evidence. The parameter
\(\gamma_c=f(0,c,c)\) is the predictive probability of a new kind after
\(c\) different successes; as \(\gamma_c\to0\) the posterior of \(C^c\) on
such evidence tends to one, and \(P(C^c)\) decreases as \(\gamma_c\)
increases, so \(\gamma_w\) is an index of caution for constituents of
width \(w\). Where the two-dimensional system has one index \(\alpha\) of
overall pessimism about constituents \(C^w\), \(w<K\), the
\(K\)-dimensional system has one for each \(w<K\). It also allows other
prior distributions, for example ones violating (19): Carnap (1968)
proposed dividing probability equally among the
\emph{constituent-structures} \(S^w\) (\(w=0,\ldots,K\)), each the
disjunction of the \(\binom Kw\) constituents of width \(w\)
\cite[pp.~324--325]{Niiniluoto2011HintikkaProgramReading}.

If the \(\gamma_c\) do not have their Carnapian values, Niiniluoto states
that one can show
\[
 (27)\qquad P(Q_i(a_{n+1})/e\&C^w)=\frac{n_i+\lambda/K}{n+w\lambda/K}.
\]
Comparing with (15), the intersection of the \(K\)-dimensional system
with the \(\lambda\)--\(\alpha\) system contains, besides the
\(\lambda\)-continuum, the members in which \(\lambda\), as a function of
\(w\), equals a constant multiple of \(w\) (printed ``\(\alpha w\) for some
constant \(a>0\)''); the multiplier \(1\) is the generalized combined
system (15\('\)). A footnote reports that Hilpinen (1968, p.~65) mentions
this class and that Kuipers calls them SH-systems; Kuipers (1978) studies
the relations among systems in detail. From (27) the \(K\)-dimensional
system satisfies Reichenbach's Axiom (3) and Positive Instantial
Relevance (4), and the adequacy condition (13) holds whenever the
parameters are more optimistic than Carnap's:
\[
 (28)\qquad \gamma_i<\delta_i\ (i=c,\ldots,K-1)\implies P(C^c/e)\to1
 \quad(n\to\infty,\ c\text{ fixed}).
\]
So the zero confirmation of laws in Carnap's continuum is an accidental
feature, removed by weakening A5 to A6. A further footnote notes that
Zabell (1997) developed Johnson's axiomatization so that constituents of
width one get nonzero probability, a special case of the
\(K\)-dimensional system with \(\gamma_1<\delta_1\)
\cite[p.~325 and nn.~3--4]{Niiniluoto2011HintikkaProgramReading}. Claims
(27) and (28), and the derivation of (3) and (4) from (27), are reported
from the chapter without proof.

\begin{hbtheorem}{thm:p15-w02b-p15-030-intersection}{proved=annex:p15-w02b-p15-030-intersection}{When the two-dimensional likelihoods coincide with (27)}
Fix \(w\geq1\), \(\lambda>0\) and \(\lambda_H>0\). The identity
\[
 \frac{n_i+\lambda_H/w}{n+\lambda_H}=\frac{n_i+\lambda/K}{n+w\lambda/K}
\]
holds for all integers \(0\leq n_i\leq n\) if and only if
\(\lambda_H=w\cdot\lambda/K\). Hence a member of the two-dimensional system
with width-dependent parameter \(\lambda_H(w)\) has the likelihoods (27)
for every \(w\) exactly when \(\lambda_H(w)=aw\) with the constant
\(a=\lambda/K>0\); \(a=1\) gives (15\('\)).
\end{hbtheorem}

\subsection{Extensions of Hintikka's system}
\label{sec:p15-w02b-p15-030-nn-extensions}

The two-dimensional continuum appeared in \emph{Aspects of Inductive
Logic} (1966), edited by Hintikka and Suppes, from an International
Symposium on Confirmation and Induction held in Helsinki in late
September 1965, continuing a seminar at Stanford that spring; Hintikka
then held a joint appointment at Helsinki and Stanford. Besides essays on
the paradoxes of confirmation (Black, Suppes, von Wright) the volume has
essays on induction by Hintikka's Finnish students Hilpinen, Tuomela and
Pietarinen \cite[pp.~325--326]{Niiniluoto2011HintikkaProgramReading}.

Minor technical extensions allow \(Q\)-predicates of different widths
(as in (2\('\))), families of predicates giving discrete quantitative
descriptions, and countably many predicates (Kuipers 1978b). The harder
extensions concern stronger languages \cite[p.~326]{Niiniluoto2011HintikkaProgramReading}.
\begin{itemize}
\item \emph{Identity} (Hilpinen 1966). With identity one can record that
  the sampled individuals are distinct (sampling without replacement).
  Numerical quantifiers need \(d\) interrelated quantifiers; the maximal
  nesting of quantifiers is the \emph{quantificational depth}, and
  constituents of depth \(d\) say of each \(Q\)-predicate that it has
  \(0,1,\ldots,d-1\), or at least \(d\), instances. A depth-\(d\)
  constituent splits into subordinate constituents at depth \(d+1\). For
  finite universes, monadic constituents with identity are equivalent to
  Carnap's structure descriptions without individual constants. Extending
  the Jerusalem system (equal weight to depth-\(d\) constituents, then to
  state descriptions), Hilpinen showed it unreasonable to project more
  than \(d\) individuals into unobserved cells or to project observed
  singularities as real ones; but all constituents positing unobserved
  singularities receive equal posterior, equal to that of the constituent
  denying them, which is unintuitive. Distributing probability evenly
  first at depth 1, then to subordinate constituents at depth 2, and so on,
  gives the highest probability to the constituent denying unobserved
  singularities \cite[p.~326]{Niiniluoto2011HintikkaProgramReading}.
\item \emph{Order} (Tuomela 1966). For the decidable language with
  \(Rxy\) (``\(y\) is an immediate successor of \(x\)''), \(Q\)-predicates
  specify triples (predecessor, object, successor) and constituents say
  which triples exist; with equal priors the simplest constituent
  compatible with evidence has the greatest posterior, so (13) is
  achievable in an ordered universe \cite[pp.~326--327]{Niiniluoto2011HintikkaProgramReading}.
\item \emph{Full first-order logic} (Hilpinen 1971). Hintikka's 1953
  distributive normal forms exist for languages with finitely many
  polyadic relations: at each depth \(d\) a formula is a disjunction of
  depth-\(d\) constituents, expandable to greater depth. Undecidability
  means some constituents are non-trivially inconsistent with no
  effective way to locate them. The \(Q\)-predicates become attributive
  constituents, trees with branches of length \(d\): a depth-1
  constituent tells what kinds of individuals exist, described by their
  properties and relations to themselves; a depth-2 constituent describes
  all kinds of pairs; a depth-\(d\) constituent \(C^{(d)}\) is a finite set
  of finite trees whose branches correspond to sequences of individuals
  drawn with replacement. These are the strongest generalizations of
  depth \(\leq d\), and every complete theory is axiomatized by a monotone
  sequence \(\langle C^{(d)}\mid d<\infty\rangle\) with
  \(\ldots C^{(d+1)}\vdash C^{(d)}\vdash\ldots\vdash C^{(1)}\). Wherever
  analogues of A1, A2, A3 and A6 hold, one depth-\(d\) constituent gets
  asymptotic probability one on evidence of ramified sequences of \(d\)
  related individuals; the general case has not been studied
  \cite[p.~327]{Niiniluoto2011HintikkaProgramReading}.
\end{itemize}
In footnotes Niiniluoto reports the probability-logic tradition (Scott
and Krauss 1966; Fenstad 1968) and its main representation theorem,
attributed to Jerzy Łoś in 1963: the probability of a quantified sentence
is a weighted average involving a measure over \(L\)-structures and
measures over sets of individuals within each structure; exchangeability
again yields instantial relevance principles (Nix and Paris 2007), but
without new fruitful applications to induction. He criticizes Nix and
Paris (2007) on binary inductive logic for ignoring Hintikka's program and
Hilpinen (1971), and calls their proposal to reduce binary relations such
as ``\(x\) pollinates \(y\)'' to unary predicates entirely implausible as a
general move: reducing relations to monadic properties was a dogma of
classical logic until De Morgan and Peirce began the serious study of
the logic of relations in the mid-nineteenth century (Kneale and Kneale
1962), and the 1930s metalogical
results (monadic logic decidable, polyadic undecidable) made the
difference crucial, with dramatic consequences for distributive normal
forms as well \cite[pp.~327--328 and nn.~5--6]{Niiniluoto2011HintikkaProgramReading}.

Hilpinen's depth-2 analysis takes evidence on \(n\) individuals with a
complete description of their pairwise relations. Constituents \(C^w\) of
depth 2 say what kinds of individuals exist; a statement \(D^v\) says, for
each observed \(a_i\), which attributive constituent it satisfies, that is,
how observed individuals are related to unobserved ones. Priors are equal
over \(C^w\); \(P(D^v/C^w)\) and \(P(e/D^v\&C^w)\) use representative
functions of the form (15), the latter applied to pairs. On large
evidence the highest posterior goes to the simplest conjunction
\(D^e\&C^e\), where \(D^e\) says observed individuals relate to unobserved
ones as to observed ones and \(C^e\) says only the kinds already
exemplified exist. Hilpinen notes a second kind of simplicity: assuming
all observed individuals bear the same relations to observed and some
unobserved individuals reduces the number of kinds; with that statement
\(D^1\) and constituent \(C^1\), \(P(D^v/C^v)\) is maximized by \(D^1\&C^1\).
Polyadic inductive methods therefore need at least two parameters
weighting the two simplicities \cite[p.~328]{Niiniluoto2011HintikkaProgramReading}.

\paragraph{Modalities.} Soshichi Uchii (1972, 1973, 1977) extended the
system to a monadic language \(L(\Box)\) with nomic necessity \(\Box\) and
possibility \(\Diamond=\sim\Box\sim\), necessity obeying S5, relevant if
lawlike generalizations need intensional notions. Nomic constituents are
\[
 (29)\qquad \bigwedge_{i\in CT}\Diamond(\exists x)Q_i(x)\ \&\ \Box(x)\Bigl[\bigvee_{i\in CT}Q_i(x)\Bigr],
\]
which Uchii calls ``a non-paradoxical causal law'', saying which kinds are
physically possible and which impossible; stronger statements are
\[
 (30)\qquad \bigwedge_{i\in H}\Diamond C_i\ \&\ \Box\Bigl[\bigvee_{i\in H}C_i\Bigr],
\]
with \(C_i\) the ordinary constituents. Such laws are Mill's ``laws of
coexistence''; ``laws of succession'' need temporal notions. Let \(B_i\) be
the nomic constituent with the same positive \(Q\)-predicates as \(C_i\).
Since actuality entails possibility, the chapter counts \(K-w\) nomic
constituents compatible with an ordinary constituent \(C^w\) of width \(w\).
Uchii's treatment makes two assumptions: \(P(C_i)=P(B_i)\) for all \(i\), and
learning the nomic constituent \(B_i\) in addition to the actual
constitution \(C_i\) does not change the probability of the evidence,
\(P(e/C_i)=P(e/C_i\&B_i)\). From (13) it follows that
\[
 (31)\qquad P(B^c/e^c_n)\to1\ (n\to\infty,\ c\text{ fixed})\iff P(B^c/C^c)=1.
\]
So asymptotic certainty that \(C^c\) truly describes the actual constitution
of the universe carries over to the nomic constituent \(B^c\) exactly when
\(P(C^c/B^c)=P(B^c/C^c)=1\), a condition on both directions of conditioning
\cite[p.~329]{Niiniluoto2011HintikkaProgramReading}.
Niiniluoto calls this a very strong assumption, and says why: it simply
eliminates all the other nomic constituents that are compatible with
\(C^c\), constituents which the asymptotic evidence \(e^c_n\) undermines but
cannot refute, since each of them allows the actual world to exhibit
exactly the \(c\) observed kinds. Uchii secures it by postulating
\(P((\exists x)\varphi(x)/\Diamond(\exists x)\varphi(x))=1\) for all formulas
\(\varphi\), which is the Principle of Plenitude: the questionable
metaphysical doctrine that every genuine possibility is realized in the
actual history \cite[p.~329]{Niiniluoto2011HintikkaProgramReading}.
\Cref{thm:p15-w02b-p15-030-uchii} proves (31) from the two assumptions, shows
that the two-sided condition reduces to one side under \(P(C^c)=P(B^c)\),
makes the elimination precise, and shows that Plenitude delivers the
condition.

Niiniluoto's alternative (1987, pp.~101--102) starts from a different
estimate of what is worth knowing. How the universe happens actually to be
constituted may matter less than it seems, because evidence gathered by
active experimentation itself brings new possibilities into being; and
since laws of nature support counterfactuals, one may hold that experiment
is what chiefly confirms them (von Wright 1971). The true actual constituent \(C^c\)
fluctuates as experiments bring new kinds into existence, while the true
nomic constituent expresses permanent features of the universe; interest
should therefore attach to the latter. The proposal is to formulate the
inductive approach of the axiomatic systems
(\cref{sec:p15-w02b-p15-030-nn-axioms}) directly with nomic constituents, so
that the axioms themselves imply convergence for \(B^c\) on experimental
evidence \(e^c_n\), with no plenitude postulate linking actual and nomic
constituents \cite[p.~329]{Niiniluoto2011HintikkaProgramReading}. The chapter
does not carry this formulation out, and none is given here.

\begin{hbtheorem}{thm:p15-w02b-p15-030-uchii}{proved=annex:p15-w02b-p15-030-uchii}{Uchii's convergence condition (31) and the Plenitude postulate}
Let \(P\) be a probability on the sentences of \(L(\Box)\) that gives
sentences equivalent in monadic S5 equal probability. For \(CT\subseteq\{1,\ldots,K\}\)
write \(C_{CT}\) for the ordinary constituent (10) and \(B_{CT}\) for the
nomic constituent (29) with positive set \(CT\); let \(CT_c\) be the set of
the \(c\) kinds observed in \(e^c_n\), \(C^c=C_{CT_c}\), \(B^c=B_{CT_c}\), and
assume \(P(C^c\&e^c_n)>0\) for every \(n\).
\begin{enumerate}
\item If \(P(e^c_n/C^c)=P(e^c_n/C^c\&B^c)\) whenever \(P(C^c\&B^c)>0\), then
  \(P(B^c\&C^c/e^c_n)=P(B^c/C^c)\,P(C^c/e^c_n)\). If moreover (13) holds,
  \(P(B^c/e^c_n)\to P(B^c/C^c)\) as \(n\to\infty\) with \(c\) fixed; in
  particular (31) holds.
\item If \(P(C^c)=P(B^c)\), then \(P(B^c/C^c)=P(C^c/B^c)\); so
  \(P(B^c/C^c)=1\) iff \(P(C^c/B^c)=P(B^c/C^c)=1\).
\item \(B_{CT}\) and \(B_{CT'}\) are jointly inconsistent when \(CT\neq CT'\),
  and \(B_{CT}\) is inconsistent with \(C_{CT_0}\) unless \(CT_0\subseteq CT\).
  In S5 Kripke models with a common domain, \(B_{CT}\) and \(C_{CT_0}\) are
  jointly satisfiable whenever \(CT_0\subseteq CT\); hence the nomic
  constituents compatible with \(C^w\) are the \(2^{K-w}\) constituents
  \(B_{CT}\) with \(CT\supseteq CT_w\).
\item If \(P(B^c/C^c)=1\), then \(P(B_{CT}/C^c)=0\) for every \(CT\neq CT_c\),
  in particular for each of the other nomic constituents compatible with
  \(C^c\).
\item If \(P\) satisfies the Plenitude postulate
  \(P((\exists x)Q_i(x)/\Diamond(\exists x)Q_i(x))=1\) for every \(i\) with
  \(P(\Diamond(\exists x)Q_i(x))>0\), then \(P(B^c/C^c)=1\).
\end{enumerate}
\end{hbtheorem}

Part~1 is the content of the chapter's ``it follows from (13)''; the
proof uses Uchii's second assumption only at \(i=c\) and needs his first
assumption only for part~2, which is why the chapter can state the
condition two-sidedly. Part~3 shows that the chapter's count \(K-w\) does
not match its own definitions: the nomic constituents compatible with
\(C^w\) are indexed by the subsets of the \(K-w\) unexemplified kinds, so
there are \(2^{K-w}\) of them, \(2^{K-w}-1\) besides \(B_{CT_w}\). The printed
count is kept above as the chapter's wording; the qualitative point, that
Uchii's condition removes every rival nomic constituent compatible with
\(C^c\), is part~4 and does not depend on the count. Part~5 uses Plenitude
only for the atomic formulas \(Q_i\); the chapter's postulate, for all
\(\varphi\), includes these instances.

\subsection{Semantic information}
\label{sec:p15-w02b-p15-030-nn-information}

Hintikka (1968) observed that his treatment of induction is not purely
enumerative: a generalization's probability depends also on the
evidence's power to refute universal statements, an eliminative aspect
related to Popper's falsification. Popper wanted bold, informative and
hence improbable theories. With prior (17), the initially least probable
compatible constituent \(C^c\) ends with the highest posterior, and the
smaller \(\alpha\), the faster confirmation shifts to initially less
probable constituents; choosing a small \(\alpha\) is, in Hintikka's
words, an indication of one aspect of the boldness Popper advocated (1966,
p.~131, as cited). The theory of semantic information (Hintikka and
Pietarinen 1966) gives a systematic argument: information content is
inversely related to prior probability, a strong generalization is more
informative the fewer kinds it admits, so \(C^c\) is the most informative
compatible constituent, and for constituents high information, low prior
and high posterior go together \cite[pp.~329--330]{Niiniluoto2011HintikkaProgramReading}.

Carnap and Bar-Hillel (1952) defined the information content of \(h\) as
the class of content elements (negations of state descriptions) it
entails, equivalently the range \(R(\sim h)\); if Popper's basic sentences
are state descriptions, this is Popper's 1934 empirical content.

\begin{hbdefinition}{def:p15-w02b-p15-030-nn-information}{Information measures (32)--(37)}
\[
\begin{aligned}
 (32)\quad&\operatorname{inf}(h)=-\log P(h),\\
 (33)\quad&\operatorname{cont}(h)=P(\sim h)=1-P(h),\\
 (34)\quad&\operatorname{inf}_{add}(h/e)=\operatorname{inf}(h\&e)-\operatorname{inf}(e),\\
          &\operatorname{cont}_{add}(h/e)=\operatorname{cont}(h\&e)-\operatorname{cont}(e),\\
 (35)\quad&\operatorname{inf}_{cond}(h/e)=-\log P(h/e),\\
          &\operatorname{cont}_{cond}(h/e)=1-P(h/e),\\
 (36)\quad&\operatorname{transinf}(h/e)=\operatorname{inf}(h)-\operatorname{inf}(h/e),\\
          &\operatorname{transcont}_{add}(h/e)=\operatorname{cont}(h)-\operatorname{cont}_{add}(h/e),\\
          &\operatorname{transcont}_{cond}(h/e)=\operatorname{cont}(h)-\operatorname{cont}_{cond}(h/e),\\
 (37)\quad&\operatorname{expl}_1(h,e)=\operatorname{transinf}(e/h)/\operatorname{inf}(e),\\
          &\operatorname{expl}_2(h,e)=\operatorname{transcont}_{add}(e/h)/\operatorname{cont}(e),\\
          &\operatorname{expl}_3(h,e)=\operatorname{transcont}_{cond}(e/h)/\operatorname{cont}(e).
\end{aligned}
\]
(32) is the surprise value, formally like Shannon's measure; (33) is
substantial information; (34) incremental, (35) conditional and (36)
transmitted information (Hintikka 1968); (37) are normalized measures of
explanatory or systematic power (Hintikka 1968; Pietarinen 1970;
Niiniluoto and Tuomela 1973), \(\operatorname{expl}_2\) being Hempel and
Oppenheim's 1948 measure. In (36), \(\operatorname{inf}(h/e)\) is
\(\operatorname{inf}_{cond}(h/e)\)
\cite[pp.~330--331]{Niiniluoto2011HintikkaProgramReading}.
\end{hbdefinition}

\begin{hbtheorem}{thm:p15-w02b-p15-030-information}{proved=annex:p15-w02b-p15-030-information}{Identities among the information measures}
Let \(P(h\&e)>0\) and \(0<P(e)<1\). Then
\begin{enumerate}
\item \(\operatorname{inf}_{add}(h/e)=\operatorname{inf}_{cond}(h/e)\);
\item \(\operatorname{transinf}(h/e)=\log P(h/e)-\log P(h)\),
  \(\operatorname{transcont}_{add}(h/e)=1-P(h\lor e)\),
  \(\operatorname{transcont}_{cond}(h/e)=P(h/e)-P(h)\);
\item \(\operatorname{expl}_1(h,e)=\dfrac{\log P(e)-\log P(e/h)}{\log P(e)}\),
  \(\operatorname{expl}_2(h,e)=\dfrac{1-P(h\lor e)}{1-P(e)}=P(\sim h/\sim e)\),
  \(\operatorname{expl}_3(h,e)=\dfrac{P(e/h)-P(e)}{1-P(e)}\);
\item \(\operatorname{transinf}(h/e)>0\iff\operatorname{transcont}_{cond}(h/e)>0\iff P(h/e)>P(h)\),
  and \(\operatorname{transcont}_{add}(h/e)>0\iff P(h\lor e)<1\);
\item if \(h\) entails \(e\), all three measures in (37) equal \(1\).
\end{enumerate}
\end{hbtheorem}

Part~5 is Niiniluoto's remark that these measures cannot discriminate
among rival deductive explanations; with inductive explanation explicated
by positive relevance they can compare rival inductive explanations
(Niiniluoto and Tuomela 1973, with a footnote on measures relative to sets
of competing hypotheses). Hilpinen (1970) used the measures to give an
information-theoretic justification of the principle of total evidence.
As both Carnap and Popper gave universal generalizations probability
zero, neither could use \(\operatorname{cont}\) to compare rival laws;
Hintikka's account opened that use, developed in ``The Varieties of
Information and Scientific Explanation'' (1968) and \emph{Information and
Inference} (1970) \cite[pp.~330--331 and n.~7]{Niiniluoto2011HintikkaProgramReading}.

\subsection{Confirmation and acceptance}
\label{sec:p15-w02b-p15-030-nn-confirmation}

Since universal hypotheses can be confirmed by finite evidence, the two
senses of confirmation separate (Carnap 1962; Niiniluoto 1972)
\cite[p.~332]{Niiniluoto2011HintikkaProgramReading}:
\begin{description}
\item[(HP)] \emph{High Probability}: \(e\) confirms \(h\) if
  \(P(h/e)>q\geq\tfrac12\);
\item[(PR)] \emph{Positive Relevance}: \(e\) confirms \(h\) iff
  \(P(h/e)>P(h)\).
\end{description}
PR is equivalent to \(P(h\&e)>P(h)P(e)\), to \(P(h/e)>P(h/\sim e)\) and to
\(P(e/h)>P(e)\) (\cref{thm:p15-w02b-p15-030-relevance}). HP satisfies
\emph{Special Consequence} (SC: if \(e\) confirms \(h\) and \(h\vdash g\),
\(e\) confirms \(g\)); PR satisfies \emph{Converse Entailment} (CE: if a
consistent \(h\) entails a non-tautological \(e\), \(e\) confirms \(h\)).

Peirce called hypothetical inference to an explanation
\emph{abduction}. Following him, Niiniluoto uses \emph{abductive
confirmation} for the support a theory receives from its
explanatory successes; the \emph{abductive criterion} is the requirement
(38) that such successes confirm. By Bayes's theorem PR meets the
abductive criterion, given \(P(h)>0\) and \(P(e)<1\):
\begin{description}
\item[(38)] if \(h\) deductively or inductively explains or predicts \(e\),
  then \(e\) confirms \(h\)
\end{description}
(Niiniluoto 1999). For deductive explanation or prediction, \(h\vdash e\),
this is part~2 of \cref{thm:p15-w02b-p15-030-relevance}; for inductive
explanation explicated by positive relevance, \(P(e/h)>P(e)\), it is the
last equivalence of part~1 of the same theorem. Further, no reasonable
notion of confirmation satisfies SC and CE together. Carnap (1962)
defined
\[
 (39)\qquad \operatorname{conf}(h,e)=P(h/e)-P(h),
\]
which is the transmitted information (36); many measures of confirmation,
support and corroboration, including Popper's (1959, p.~400, as cited),
are variants of it. Popper was right that corroboration is neither prior
nor posterior probability, but in Hintikka's system, with enough
evidence, \(C^c\) maximizes both posterior probability and
\(\operatorname{cont}\), hence (39), written \(P(h/e)+\operatorname{cont}(h)-1\)
\cite[p.~332]{Niiniluoto2011HintikkaProgramReading}.

Hintikka (1968) proposed a new measure of corroboration whose stated
merit is an interesting treatment of \emph{weak} generalizations, that
is, generalizations other than constituents, which are the strong ones
(p.~319); on this treatment see also Hintikka and Pietarinen (1966) and
Niiniluoto and Tuomela (1973). For \(h\) equivalent to
\(C_1\lor\ldots\lor C_m\) it is the least posterior probability of a
disjunct,
\[
 (40)\qquad \operatorname{corr}(h,e)=\min\{P(C_i/e)\mid i=1,\ldots,m\},
\]
which, unlike probability, covaries with logical strength:
\[
\begin{aligned}
 (41)&\quad e\vdash h_1\supset h_2\implies P(h_1/e)\leq P(h_2/e);\\
 (42)&\quad e\vdash h_1\supset h_2\implies \operatorname{corr}(h_1,e)\geq\operatorname{corr}(h_2,e);
\end{aligned}
\]
and
\[
 (43)\quad\text{with large enough }e^c_n\ (c\text{ fixed}),\ \operatorname{corr}(h,e^c_n)
 \text{ is maximal when }h\text{ is }C^c.
\]
\cite[pp.~332--333]{Niiniluoto2011HintikkaProgramReading}

\begin{hbtheorem}{thm:p15-w02b-p15-030-relevance}{proved=annex:p15-w02b-p15-030-relevance}{Positive relevance, the principal monotonicities and their hypotheses}
Let \(0<P(e)<1\) and \(P(h)>0\).
\begin{enumerate}
\item \(P(h/e)>P(h)\iff P(h\&e)>P(h)P(e)\iff P(h/e)>P(h/\sim e)\iff
  P(e/h)>P(e)\).
\item (CE) If \(h\vdash e\) and \(P(e)<1\), then \(P(h/e)>P(h)\).
\item (41) holds for every probability.
\item Let \(h_1,h_2\) be consistent generalizations with normal forms
  (11), and let \(\operatorname{corr}^e(h,e)\) be the minimum of
  \(P(C_i/e)\) over the disjuncts \(C_i\) of \(h\) that are consistent with
  \(e\). If \(e\vdash h_1\supset h_2\), then
  \(\operatorname{corr}^e(h_1,e)\geq\operatorname{corr}^e(h_2,e)\). If
  \(\vdash h_1\supset h_2\), then also
  \(\operatorname{corr}(h_1,e)\geq\operatorname{corr}(h_2,e)\) for the
  minimum (40) over all disjuncts.
\item With the minimum over all disjuncts, (42) can fail under the
  hypothesis \(e\vdash h_1\supset h_2\) alone.
\end{enumerate}
\end{hbtheorem}

Part~5 shows that the printed hypothesis of (42) requires the minimum in
(40) to range over the constituents compatible with the evidence, or the
hypothesis to be strengthened to \(\vdash h_1\supset h_2\).

\paragraph{Paradoxes.} In Hempel's paradox ``All ravens are black'' is
confirmed by black ravens, black non-ravens and non-black non-ravens. The
standard Bayesian answer, due to Hosiasson-Lindenbaum in 1940, is that
these give different incremental confirmation because in a finite universe
the cells have different sizes. Instead of that empirical assumption, one
may stipulate different widths for ``black'', ``non-black'', ``raven'' and
``non-raven'' and use (2\('\)); Hintikka (1969) motivates this ``inductive
asymmetry'' by an ordering of the primitive predicates. For Goodman's
``grue'', Hintikka's solution uses \(\alpha\) as a measure of the
confirmability of laws: assign a smaller \(\alpha\) to the scheme with
``green'' than to the one with ``grue'', and the different confirmation of
``All emeralds are green'' and ``All emeralds are grue'' is explained
\cite[p.~333]{Niiniluoto2011HintikkaProgramReading}.

\paragraph{Acceptance.} Carnap's system has no acceptance rules: the
task of inductive logic is to evaluate the epistemic probabilities of
hypotheses, which are then used in decision making (Carnap 1980;
Stegmüller 1973). Here Carnap agrees with frequentist statisticians
(Neyman, E.~S. Pearson) and with the Bayesian Savage, who recommend that
inductive inference give way to inductive behaviour, that is, to actions
based on probabilities. On this view the scientist's main role is to
advise practical decision makers, not to seek new truths. On the
cognitivist model of inquiry, by contrast, the tentative results of
research make up a body of accepted hypotheses, the ``scientific
knowledge'' of a given time. In the spirit of Peirce's fallibilism these
hypotheses remain open to question and revision at any time, by new
evidence or new theoretical insight; yet under some conditions it is
rational to accept a hypothesis tentatively on the strength of evidence,
and one task of inductive logic is to define acceptance rules for such
corrigible factual statements. Hintikka, together with Levi (1967), is a
cognitivist \cite[pp.~333--334]{Niiniluoto2011HintikkaProgramReading}. With the accepted set
consistent and closed under consequence, Kyburg's lottery paradox shows
high probability is not enough. In Hintikka's system a threshold sample
size \(n_0\) can be computed so that \(C^c\) exceeds \(1-\varepsilon\):
\[
 (44)\quad\text{if }P(C^c/e)\geq1-\varepsilon\iff n\geq n_0,\text{ then on }
 e\text{ accept }C^c\iff n\geq n_0
\]
(Hintikka and Hilpinen 1966; Hilpinen 1968), and by closure all
generalizations entailed by \(C^c\). In the two-dimensional continuum
\(n_0\) is the largest integer \(n\) for which
\[
 \varepsilon'\leq\max\sum_{i=1}^{K-c}\binom{K-c}{i}\Bigl(\frac c{c+i}\Bigr)^{n-\alpha},
\]
the maximum taken over \(0\leq c\leq K-1\), with
\(\varepsilon'=\varepsilon/(1-\varepsilon)\). Hintikka and Hilpinen also
hold
\[
 (45)\quad \varphi(a_i)\text{ is acceptable on }e\iff(x)\varphi(x)
 \text{ is acceptable on }e,
\]
reducing singular inference (Mill's ``eduction'') to universal inference.
A footnote distinguishes (44), a factual detachment rule concluding a
factual statement from factual and probabilistic premises, from
probabilistic detachment rules that are deductive arguments within the
calculus (Suppes 1966), for example from \(P(h/e)=r\) and \(P(e)=1\) to
\(P(h)=r\) \cite[pp.~333--334 and n.~8]{Niiniluoto2011HintikkaProgramReading}.
The threshold formula is reported as printed, without derivation; at
\(c=0\) its summand \((c/(c+i))^{n-\alpha}\) is \(0\).

\subsection{Cognitive decision theory}
\label{sec:p15-w02b-p15-030-nn-decision}

Bayesian decision theory recommends the act maximizing subjective
expected utility. Cognitive decision theory, anticipated by Bolzano and
proposed independently by Hempel and Levi in the early 1960s, applies the
same rule to acceptance of hypotheses with epistemic utilities (truth,
information, explanatory and predictive power, simplicity); inductive
logic may supply the probabilities
\cite[pp.~334--335]{Niiniluoto2011HintikkaProgramReading}.

\begin{hbdefinition}{def:p15-w02b-p15-030-nn-problem}{Cognitive problem, answers and expected utility}
Let \(B=\{h_1,\ldots,h_n\}\) be mutually exclusive and jointly exhaustive
hypotheses (for example state descriptions, structure descriptions or
constituents of a monadic language, or complete theories of a finite
first-order language); interpreted on a domain \(U\), exactly one, \(h^*\),
is true, the \emph{target}. The \emph{complete answers} are the \(h_i\);
the \emph{partial answers} \(D(B)\) are the nonempty disjunctions of
complete answers, the tautology being ``I don't know''. With epistemic
utility \(u(g,h_j)\) of accepting \(g\in D(B)\) when \(h_j\) is true and a
rational probability \(P\),
\[
 (46)\qquad U(g/e)=\sum_{j=1}^nP(h_j/e)\,u(g,h_j),
\]
with the rules (EU) accept the \(g\in D(B)\) maximizing \(U(g/e)\), and
(CP) moving from \(g\) to \(g'\) is cognitively progressive on \(e\) iff
\(U(g/e)<U(g'/e)\) (Niiniluoto 1995). For \(g\) with
\((47)\ \vdash g\equiv\bigvee_{i\in I_g}h_i\), the truth-value utility is
\(u(g,h_j)=1\) if \(h_j\) is in \(g\) and \(0\) otherwise, so \(u(g,h^*)\) is
the truth value \(\operatorname{tv}(g)\). A footnote notes that \(B\) may be a
subset of a quantitative space such as \(\mathbb R\), sums becoming integrals
\cite[pp.~335--336 and n.~9]{Niiniluoto2011HintikkaProgramReading}.
\end{hbdefinition}

The candidate utilities are (49) Levi's \(\operatorname{tv}(g)+qI(g)\), with
\(I(g)\) the normalized number of complete answers excluded by \(g\) and
\(0<q\leq1\) an index of boldness; (51) \(\operatorname{cont}(g)\) if \(g\) is
true and \(-\operatorname{cont}(\sim g)\) if false (Levi 1963, rejected in
1967, defended by Hintikka and Pietarinen 1966); (53) \(P(\sim g/\sim e)\) if
true and \(-P(g/\sim e)\) if false (Hempel's systematic power; Pietarinen
1970); and (54) Hilpinen's (1968) \(1-P(g)\) if true and \(-qP(g)\) if false
\cite[pp.~336--337]{Niiniluoto2011HintikkaProgramReading}.

\begin{hbtheorem}{thm:p15-w02b-p15-030-utilities}{proved=annex:p15-w02b-p15-030-utilities}{Expected epistemic utilities (48), (50), (52), (53), (55)}
Let \(g\in D(B)\) and let \(P(\cdot/e)\) be a probability on \(B\).
\begin{enumerate}
\item With the truth-value utility, \(U(g/e)=P(g/e)\) (48). With
  \(u=\operatorname{cont}(g)\) regardless of truth, \(U(g/e)=1-P(g)\).
\item With Levi's utility (49), \(U(g/e)=P(g/e)+qI(g)\) (50).
\item With (51), \(U(g/e)=P(g/e)-P(g)\) (52).
\item With (53) and \(0<P(e)<1\),
  \(U(g/e)=P(g/e)-P(g/\sim e)=[P(g/e)-P(g)]/P(\sim e)\).
\item With (54) as printed,
  \(U(g/e)=P(g/e)-qP(g)-(1-q)P(g)P(g/e)\), which equals (55)
  \(P(g/e)-qP(g)\) iff \(q=1\) or \(P(g)P(g/e)=0\). With the true-case
  utility \(1-qP(g)\) in place of \(1-P(g)\), \(U(g/e)=P(g/e)-qP(g)\)
  exactly.
\end{enumerate}
\end{hbtheorem}

Niiniluoto's consequences: with the truth-value utility posterior
probability is expected truth value, and maximizing it is extremely
conservative, since the best answers are those certain on \(e\), such as
\(e\) itself or a tautology, and weakening an uncertain answer is always
progressive; Popper made this argument against probability as a criterion
in 1934. With information alone, a contradiction of probability zero is
best and any false theory improves by adding falsehoods; the same holds
for explanatory power alone. Levi's (50) balances the two goals and makes
all partial answers comparable; (52) combines boldness and posterior
probability. In Hintikka's system, on large enough evidence, (50), (52)
and (55) are maximized by \(C^c\); since \(C^c\) is also the simplest
compatible constituent, Hintikka's posteriors favour simplicity without
adding it as an extra condition
\cite[pp.~336--337]{Niiniluoto2011HintikkaProgramReading}. Part~5 of the
theorem records that the printed (55) follows from a true-case utility
\(1-qP(g)\); with the printed \(1-P(g)\) it holds only for \(q=1\).

\subsection{Inductive logic and theories}
\label{sec:p15-w02b-p15-030-nn-theories}

Inductive logic has a reputation as a tool of narrow empiricism. Aristotle
tied induction to concept formation; from Bacon to Whewell, Mill and Jevons
and the Cambridge school its role in science was stressed, and many
empiricists were inductivists. The hypothetico-deductive method lets
theories be invented freely and tested by their consequences; Bayes's
theorem evaluates such tests, and in (38) \(h\) may postulate unobservables,
so inductive probability is not committed to narrow empiricism (Hempel
1965; Niiniluoto and Tuomela 1973). Carnap's
evidence is singular observational sentences, and the zero confirmation
of laws seemed to support the ``atheoretical thesis''. Replying to Putnam
(Schilpp, p.~988), Carnap granted that inductive logic should be built for
the total language of science, observational plus theoretical, so that
predictions can take account of proposed laws, and that meaning
postulates of theories should get \(m\)-value one
\cite[pp.~337--338]{Niiniluoto2011HintikkaProgramReading}.

Hempel's ``Theoretician's Dilemma'' (1958): if a theory \(T\) in \(L\)
achieves deductive systematization between observational \(e\) and \(e'\)
in \(L_o\) (\((T\&e)\vdash e'\) but not \(e\vdash e'\)), the elimination
methods of Ramsey and Craig yield an observational subtheory achieving the
same, so theoretical terms are logically dispensable for deductive
systematization. Hempel suggested they may be indispensable for inductive
systematization, with the example
\(T=(x)(Mx\supset O_1x)\&(x)(Mx\supset O_2x)\): infer \(Ma\) from \(O_1a\)
by the first law and \(O_2a\) from \(Ma\) by the second. Niiniluoto (1972)
shows this link is unjustified, relying at once on the incompatible CE and
SC. Tuomela studied the deductive gains of theoretical terms; the
inductive gains were the topic of Niiniluoto's 1973 dissertation
\cite[pp.~338--339]{Niiniluoto2011HintikkaProgramReading}.

Writing \(eIh\) for ``\(h\) is inducible from \(e\)'', \(T\) achieves
inductive systematization between \(e\) and \(e'\) if \((T\&e)Ie'\) but not
\(eIe'\); with PR as the explicatum of \(I\),
\[
 (IS)\quad P(e'/e\&T)>P(e'),\ \ P(e'/e)=P(e');
 \qquad (IS')\quad P(e'/e\&T)>P(e'/T),
\]
both presupposing probabilities conditional on genuine theories.
Calculations in Niiniluoto and Tuomela (1973) use the generalized combined
system: for a universal generalization \(h\) in \(L_o\) and a theory \(T\)
with theoretical terms, \(P(h/e\&T)\) depends on the number \(b_0\) of
\(Q\)-predicates of \(L\) empty by \(h\) but not by \(T\), and for large \(n\)
and \(\alpha\)
\[
 (56)\qquad P(h/e\&T)\simeq\frac1{(1+\alpha/n)^{b_0}};
\]
comparison with (24) gives \(P(h/e)<P(h/e\&T)\) iff \(b>b_0\) (op.\ cit.,
p.~38). If \(T\) explicitly defines \(M\) in \(L_o\), then \(b=b_0\), so
\[
 (57)\qquad P(h/e\&T)=P(h/e).
\]
The main result of their ch.~9 is that theoretical terms can be logically
indispensable for observational inductive systematization, which helps
solve the dilemma and supports scientific realism. Probabilities
\(P(h/e\&T)\) also serve hypothetico-inductive inference with a tentative
\(T\) (observational support from \(e\), theoretical support from \(T\));
one can compute \(P(C^w/e)\) for constituents of the full language on
observational evidence (Niiniluoto 1976). These applications show that
inductive logic can be developed as a non-inductivist methodological
program (Niiniluoto and Tuomela 1973, ch.~12). Changes of \(\lambda\) and
\(\alpha\) cannot be modelled by conditionalization (compare Hintikka
1966, 1987b, 1997); Niiniluoto and
Tuomela keep parameters fixed and conditionalize on conceptual
information, and (57) gives linguistic invariance: probabilities in \(L\)
conditional on an explicit definition equal those in \(L_o\), and
intertranslatable languages have equal inductive logics (op.\ cit.,
p.~175) \cite[pp.~339--340]{Niiniluoto2011HintikkaProgramReading}. Within
the approximations (24) and (56), and since \(1+\alpha/n>1\) for
\(\alpha>0\), the comparison \(b>b_0\) is just monotonicity of
\(x\mapsto(1+\alpha/n)^{-x}\); the exact claim rests on the cited
calculations, which are not reproduced here.

\subsection{Analogy and observational errors}
\label{sec:p15-w02b-p15-030-nn-analogy}

Analogy generalizes the identity rule (RI) from \(F(a)\) and \(b=a\) to
\(F(b)\), replacing identity by similarity: (RS) from \(F(a)\) and ``\(b\) is
similar to \(a\)'' to \(F(b)\). Classically similarity is partial identity;
in Keynes's terms shared attributes form the positive analogy and others
the negative analogy; Mill measured the strength by \(k/(k+m)\) for \(k\)
shared and \(m\) differing attributes, making analogy enumerative
induction over properties. Carnap's and Hintikka's systems handle simple
positive analogy but fail with negative analogy (Hesse 1964)
\cite[p.~340]{Niiniluoto2011HintikkaProgramReading}.

\begin{hbtheorem}{thm:p15-w02b-p15-030-analogy}{proved=annex:p15-w02b-p15-030-analogy}{The analogy formulas (58)--(60)}
In a monadic language with \(k\) primitive predicates (\(K=2^k\)), let
\(F_1(x)=M_1(x)\&\ldots\&M_m(x)\) and
\(F_2(x)=M_1(x)\&\ldots\&M_n(x)\) with \(1\leq m<n\leq k\), of widths
\(w_1=2^{k-m}>w_2=2^{k-n}\). Let \(P\) satisfy A1--A4 and A6, and put
\(\varphi=f(0,1,1)\), the probability that \(b\) is of a given kind other
than that of a distinct individual \(a\) whose kind is known. Then
\[
 P(F_2(b)/F_1(b))=\frac{w_2}{w_1},\qquad
 P(F_2(b)/F_1(b)\&F_2(a))=\frac{1-(K-w_2)\varphi}{1-(K-w_1)\varphi},
\]
and the second exceeds the first if and only if \(\varphi<1/K\)
(equivalently \(f(1,1,1)>\varphi\)). If
\(\varphi=(\lambda/K)/(1+\lambda)\) with \(0<\lambda<\infty\), the second
value is \((1+w_2\lambda/K)/(1+w_1\lambda/K)\) (59), and at \(\lambda=K\) it is
\((1+w_2)/(1+w_1)\) (60).
\end{hbtheorem}

Niiniluoto prints (58) as the two values joined by ``\(>\)'' (citing his
1981 paper, p.~7) and (60) as Carnap's analogy formula for \(c^*\) (LFP,
pp.~569--570; compare \cref{thm:p15-carnap-analogy}). The condition
\(\varphi<1/K\) holds in Carnap's continuum and whenever \(\varphi=\gamma_1\)
satisfies (26) with \(1+\lambda>0\). As soon as the evidence records a known
difference between \(a\) and \(b\), he notes, the analogy influence in (58)
disappears \cite[pp.~340--341]{Niiniluoto2011HintikkaProgramReading}.

Remedies follow. Carnap's ``analogy parameter'' \(\eta\) (Carnap and
Stegmüller 1959) was applied to Hintikka's system by Pietarinen (1972).
Carnap's posthumous ``Basic System'' (1980) uses distances between
predicates: for \(Q_u,Q_v\) of the form \((\pm)M_1x\&\ldots\&(\pm)M_kx\),
\(d(Q_u,Q_v)=d_{uv}=m/k\) if they disagree on \(m\) of the \(k\) primitives
(with two primitives, \(d(M_1x\&M_2x,M_1x\&\sim M_2x)=\tfrac12\) and
\(d(M_1x\&M_2x,\sim M_1x\&\sim M_2x)=1\)); with families \(\mathbf M_j\) and
internal distances \(d_j\), a Euclidean metric over the \(k\) dimensions
makes the \(Q\)-predicates a metric space of concepts. Resemblance is
\(r_{uv}=1-d_{uv}\) (or \(1/(1+d_{uv})\)), complete similarity being
\(r_{uv}=1\). \emph{Similarity influence} modifies (4): observing \(Q_i\)
raises the expectation of nearby cells,
\[
 P(Q_2(b)/Q_1(a))>P(Q_3(b)/Q_1(a))\iff d_{12}<d_{13}\iff r_{12}>r_{13}
\]
(Carnap 1980, p.~46), which already helps with negative analogy; in
Hintikka's or the \(K\)-dimensional system \(P(C^w/e)\) should reflect
distances between \(CT_w\) and the observed cells \(CT_e\)
\cite[pp.~341--342]{Niiniluoto2011HintikkaProgramReading}.

Niiniluoto's proposals (1980, 1981) multiplied \(P(Q_i(a_{n+1})/e)\) by a
factor for the minimum distance of \(Q_i\) from \(CT_e\) or the weighted
influence of observed individuals, but did not always satisfy (4) and (3).
Kuipers (1984, 1988) lets analogy influence vanish as evidence grows,
replacing (2) by
\[
 (61)\qquad \frac{n_i+\alpha_i(e_n)+\lambda/K}{n+\alpha(n)+\lambda},
 \qquad \alpha(n)=\sum_{i=1}^K\alpha_i(e_n),
\]
with \(\alpha_i(e_n)\) the analogy profit of \(Q_i\) and the marginal analogy
\(\alpha(n)-\alpha(n-1)\) decreasing from a positive value to zero, which
secures (3); the \(K\)-dimensional likelihoods can be modified similarly
(Niiniluoto 1988). A footnote lists other treatments (Skyrms 1993, di Maio
1996, Festa 1996, Maher 2000, Romeyn 2005)
\cite[p.~342 and n.~10]{Niiniluoto2011HintikkaProgramReading}. The
\(\alpha_i\) of (61) are Kuipers's analogy profits, not Hintikka's
\(\alpha\).

For observational errors (Niiniluoto 1997): Hintikka's likelihoods (15)
restrict \(Q_i\) to cells allowed by \(C^w\), excluding error; Jeffrey
conditionalization changes probabilities directly, whereas Bayesian
statistics uses an error distribution. With \(Q_j(a)\) ``\(a\) is \(Q_j\)'' and
\(S_j(a)\) ``\(a\) seems to be \(Q_j\)'', assume
\[
 (62)\qquad p_{ij}=P(S_i(a)/Q_j(a))\text{ decreases with }d_{ij}=d(Q_i,Q_j),
\]
with constant reliability \(p_{ii}=\beta\); error-free logic has \(p_{ij}=1\)
for \(i=j\) and \(0\) otherwise. Bayes's theorem then gives the
probabilities \(P(Q_i(a)/S_j(a))\), and the probabilities
\[
 P(S_i(a_{n+1})/e_n(S)\&C^w)
\]
must be defined for evidence \(e_n(S)\) stated in \(S\)-predicates
\cite[pp.~342--343]{Niiniluoto2011HintikkaProgramReading}.

\subsection{Truthlikeness}
\label{sec:p15-w02b-p15-030-nn-truthlikeness}

In 1960 Popper defined comparative and quantitative verisimilitude, an
objective closeness to truth combining truth with information content,
with his corroboration as epistemic indicator; the definition failed
because it could not compare two false theories. From 1974 Hilpinen,
Tichý, Niiniluoto, Tuomela and Oddie used likeness between possible worlds
or Hintikkian constituents; Niiniluoto (1977) proposed estimating
verisimilitude with inductive probabilities
\cite[p.~343]{Niiniluoto2011HintikkaProgramReading}.

\begin{hbdefinition}{def:p15-w02b-p15-030-nn-tr}{Distances, approximate truth and truthlikeness (63)--(70)}
Relative to a cognitive problem \(B=\{h_i\mid i\in I\}\) with target
\(h^*\), let \(\Delta:B\times B\to\mathbb R\), \(\Delta(h_i,h_j)=\Delta_{ij}\),
with \(0\leq\Delta_{ij}\leq1\) and \(\Delta_{ij}=0\) iff \(i=j\). For monadic
constituents with classes \(CT_i\) the Clifford distance is
\[
 (63)\qquad |CT_i\,\triangle\, CT_j|/K .
\]
For \(g\in D(B)\) with \(\vdash g\equiv\bigvee_{i\in I_g}h_i\),
\[
\begin{aligned}
 (64)\quad&\Delta_{\min}(h_i,g)=\min_{j\in I_g}\Delta_{ij},\qquad
 (65)\quad AT(g,h^*)=1-\Delta_{\min}(h^*,g),\\
 (67)\quad&\Delta_{\mathrm{sum}}(h_i,g)=\sum_{j\in I_g}\Delta_{ij}\Big/\sum_{j\in I}\Delta_{ij},\\
 (68)\quad&\Delta^{\gamma\gamma'}_{ms}(h_i,g)=\gamma\,\Delta_{\min}(h_i,g)+\gamma'\,\Delta_{\mathrm{sum}}(h_i,g)\quad(\gamma>0,\ \gamma'>0),\\
 (69)\quad&Tr(g,h^*)=1-\Delta^{\gamma\gamma'}_{ms}(h^*,g),\qquad
 (70)\quad \operatorname{ver}(g/e)=\sum_{i\in I}P(h_i/e)\,Tr(g,h_i).
\end{aligned}
\]
\(g\) is approximately true if \(\Delta_{\min}(h^*,g)\) is small, and should
satisfy (66) \(Tr(g,h^*)=1\) iff \(\vdash g\equiv h^*\); \(\gamma\) and \(\gamma'\)
weigh the desire to find truth and to avoid error. \(g'\) is more truthlike
than \(g\) if \(\Delta^{\gamma\gamma'}_{ms}(h^*,g')<\Delta^{\gamma\gamma'}_{ms}(h^*,g)\),
and seems more truthlike on \(e\) if \(\operatorname{ver}(g/e)<\operatorname{ver}(g'/e)\)
\cite[pp.~343--345]{Niiniluoto2011HintikkaProgramReading}.
\end{hbdefinition}

The distance must be fixed per problem, with canonical choices: counting
syntactic differences for state descriptions, structure descriptions or
constituents; the Euclidean metric for point estimates of a real
parameter (on \(\mathbb R\) or \(\mathbb R^K\)); Minkowski metrics between
functions for quantitative laws. Comparing by both \(\Delta_{\min}\) and
\(\Delta_{\mathrm{sum}}\) would leave few answers comparable, which (68)
remedies. A practical rule: choose \(\gamma,\gamma'\) so that complete answers
nearest \(h^*\) beat ignorance \(1-\gamma'\) while the most misleading are
worse (for example \(\gamma'/\gamma\approx1/2\)). Maximizing \(\operatorname{ver}\)
is an acceptance rule comparable to Bayesian decision with loss
proportional to distance from the truth, and supports Bayesian interval
estimation (Niiniluoto 1986, 1987; Festa 1986); for monadic constituents
with Hintikka's measure, explicit formulas for \(\operatorname{ver}\) exist
(Niiniluoto 1987, ch.~9.5; 2005). With an even distribution on \(B\) the
values \(\operatorname{ver}(h_j/e)\) are all equal; and \(\operatorname{ver}(g/e)\)
can be high even when \(e\) refutes \(g\) and \(P(g/e)=0\), unlike (39) and (40)
\cite[pp.~344--345]{Niiniluoto2011HintikkaProgramReading}.

\begin{hbtheorem}{thm:p15-w02b-p15-030-truthlikeness}{proved=annex:p15-w02b-p15-030-truthlikeness}{Properties of \(Tr\) and \(\operatorname{ver}\)}
Let \(B\) be finite with \(|B|=m\geq2\), \(\Delta\) as in
\cref{def:p15-w02b-p15-030-nn-tr}, \(\gamma,\gamma'>0\).
\begin{enumerate}
\item (66) holds: \(Tr(g,h^*)=1\) iff \(I_g=\{h^*\}\).
\item For a tautology \(t\), \(Tr(t,h^*)=1-\gamma'\). For a complete answer,
  \(Tr(h_i,h^*)=1-\bigl(\gamma+\gamma'/\textstyle\sum_{j}\Delta_{*j}\bigr)\Delta_{*i}\),
  which decreases as \(\Delta(h_i,h^*)\) increases.
\item If \(\Delta_{ij}=1\) for all \(i\neq j\), then
  \(Tr(g,h^*)=\kappa+\mu\bigl[\operatorname{tv}(g)+qI(g)\bigr]\) with
  \(\mu=\gamma+\gamma'/(m-1)>0\), \(q=\gamma'/\mu\), \(\kappa\) a constant and
  \(I(g)=(m-|I_g|)/(m-1)\); here \(q\leq1\) iff \(\gamma'(m-2)\leq\gamma(m-1)\).
  So \(Tr\) is then an increasing affine function of Levi's utility (49).
\item (71): \(\sum_iP(h_i/e)\operatorname{tv}(g,h_i)=P(g/e)\). If \(P(h_j/e)=1\),
  then \(\operatorname{ver}(g/e)=Tr(g,h_j)\).
\item If \(B\) is the set of consistent monadic constituents and (13)
  holds, then \(\operatorname{ver}(g/e)\to Tr(g,C^c)\) (72), and
  \(\operatorname{ver}(g/e)\to1\) iff \(\vdash g\equiv C^c\) (73).
\end{enumerate}
\end{hbtheorem}

Part~3 makes precise Niiniluoto's remark that with the trivial distance
\(Tr\) reduces to a special case of Levi's definition; the reduction is
exact up to a positive affine rescaling and requires the stated inequality
for Levi's restriction \(q\leq1\). By (71) \(Tr:\operatorname{ver}=\operatorname{tv}:P\):
expected verisimilitude estimates real truthlikeness as posterior
probability estimates truth value, so standard Bayesianism maximizes
expected truth value while \(\operatorname{ver}\) maximizes expected
truthlikeness; by (72), asymptotically the boldest constituent compatible
with the evidence has the highest estimated verisimilitude
\cite[pp.~345--346]{Niiniluoto2011HintikkaProgramReading}.

\subsection{Machine learning}
\label{sec:p15-w02b-p15-030-nn-ml}

Other combinations of truth-closeness and probability are \emph{probable
verisimilitude} and \emph{probable approximate truth}. For \(g\in D(B)\) and
small \(\varepsilon>0\) let
\[
 (74)\qquad V_\varepsilon(g)=\{h_i\in B\mid\Delta_{\min}(h_i,g)\leq\varepsilon\},
\]
and let \(g^\varepsilon\), the ``blurred'' \(g\), be the disjunction of
\(V_\varepsilon(g)\); then \(g\vdash g^\varepsilon\), and \(g\) is approximately true
within \(\varepsilon\) iff \(g^\varepsilon\) is true. The posterior probability that
\(AT(g,h^*)\geq1-\varepsilon\) is
\[
 (75)\qquad PAT_{1-\varepsilon}(g/e)=P(h^*\in V_\varepsilon(g)/e)=\sum_{h_i\in V_\varepsilon(g)}P(h_i/e).
\]
\cite[pp.~346--347]{Niiniluoto2011HintikkaProgramReading}

\begin{hbdefinition}{def:p15-w02b-p15-030-nn-pv}{Probable verisimilitude}
Niiniluoto describes \emph{probable verisimilitude} in words only: it is
the probability, given \(e\), that \(g\) is truthlike at least to a given
degree \cite[p.~346]{Niiniluoto2011HintikkaProgramReading}. Writing the
given degree as \(1-\varepsilon\) and using (69) with the unknown target
\(h^*\) ranging over \(B\), this is
\[
 P\bigl(Tr(g,h^*)\geq1-\varepsilon\,/\,e\bigr)
 =\sum_{h_i\in B:\ Tr(g,h_i)\geq1-\varepsilon}P(h_i/e),
\]
since the event that \(Tr(g,h^*)\geq1-\varepsilon\) is the disjunction of the
complete answers \(h_i\) for which \(Tr(g,h_i)\geq1-\varepsilon\). The display is
the reading used here; the chapter prints no formula and no abbreviation
for this notion, and none is introduced. It differs from (75) only in
measuring closeness by \(Tr\) instead of \(AT\).
\end{hbdefinition}

The two notions share a feature that separates them from posterior
probability: according to Niiniluoto, the value of each can be non-zero
for a hypothesis whose probability on the evidence is zero, so that
\(PAT_{1-\varepsilon}(g/e)>0\) is possible while \(P(g/e)=0\)
\cite[p.~347]{Niiniluoto2011HintikkaProgramReading}. Part~2 of
\cref{thm:p15-w02b-p15-030-pat} proves this for both, and
\cref{def:p15-w02b-p15-030-nn-pat-worked} computes an instance.

\begin{hbdefinition}{def:p15-w02b-p15-030-nn-esc}{Evidential success condition}
\begin{description}
\item[(ESC)] Evidence \(e\) is true and fully informative about the variety
  of the world \(w\): relative to the expressive power of \(L\), it exhibits
  every kind of individual that exists in \(w\).
\end{description}
\cite[p.~347]{Niiniluoto2011HintikkaProgramReading}
\end{hbdefinition}

\begin{hbtheorem}{thm:p15-w02b-p15-030-pat}{proved=annex:p15-w02b-p15-030-pat}{Probable approximate truth and convergence under ESC}
Let \(B\) be finite and \(\Delta\) as in \cref{def:p15-w02b-p15-030-nn-tr}.
\begin{enumerate}
\item \(P(g/e)\leq PAT_{1-\varepsilon}(g/e)\), \(PAT_{1-\varepsilon}\) is
  nondecreasing in \(\varepsilon\), and at \(\varepsilon=0\),
  \(PAT_1(g/e)=P(g/e)\).
\item \(PAT_{1-\varepsilon}(g/e)>0\) iff \(P(g^\varepsilon/e)>0\). Both
  \(PAT_{1-\varepsilon}(g/e)>0\) and \(P(Tr(g,h^*)\geq1-\varepsilon\,/\,e)>0\) can
  hold while \(P(g/e)=0\).
\item If \(g_1\vdash g_2\), then \(PAT_{1-\varepsilon}(g_1/e)\leq PAT_{1-\varepsilon}(g_2/e)\).
\item Let \(B\) be the consistent monadic constituents, let \(C^*\) be the
  true one, and suppose the evidence satisfies ESC, so that the kinds
  exhibited in \(e\) are exactly those that exist. Then \(C^c=C^*\), and (13),
  (72) and (73) become (13\('\)) \(P(C^*/e)\to1\), (72\('\))
  \(\operatorname{ver}(g/e)\to Tr(g,C^*)\) and (73\('\))
  \(\operatorname{ver}(g/e)\to1\) iff \(\vdash g\equiv C^*\), as \(n\to\infty\)
  with \(c\) fixed.
\item Let \(B\) be the consistent monadic constituents and suppose (13).
  As \(n\to\infty\) with \(c\) fixed, \(PAT_{1-\varepsilon}(g/e)\to1\) if
  \(C^c\in V_\varepsilon(g)\) and \(\to0\) otherwise; and
  \(P(Tr(g,h^*)\geq1-\varepsilon\,/\,e)\to1\) if \(Tr(g,C^c)\geq1-\varepsilon\) and
  \(\to0\) otherwise.
\item If in addition ESC holds, then as \(n\to\infty\) with \(c\) fixed,
  \(PAT_{1-\varepsilon}(g/e)\to1\) iff \(g\) is approximately true within
  \(\varepsilon\), that is, iff \(\Delta_{\min}(C^*,g)\leq\varepsilon\), and
  \(PAT_{1-\varepsilon}(g/e)\to0\) otherwise; and
  \(P(Tr(g,h^*)\geq1-\varepsilon\,/\,e)\to1\) iff \(Tr(g,C^*)\geq1-\varepsilon\),
  and \(\to0\) otherwise.
\end{enumerate}
\end{hbtheorem}

After (13\('\)), (72\('\)) and (73\('\)) the chapter adds only that similar
modifications can be made in the convergence results about probable
approximate truth \cite[p.~347]{Niiniluoto2011HintikkaProgramReading}. It
states neither those results nor their modified forms; parts~5 and~6 are
the forms proved here, part~5 being the convergence of probable
approximate truth and probable verisimilitude under (13) and part~6 its
reformulation as convergence to the truth under ESC. As with (13) and
(13\('\)), part~5 says only where the degrees of belief settle, while
part~6 says that they settle on the right verdict about the actual
world.

Niiniluoto states part~2 with the unconditional \(P(g^\varepsilon)>0\); since
\(PAT\) is by (75) the conditional probability of \(g^\varepsilon\), the
conditional form is the one proved, and the two agree whenever the
posterior on \(B\) is regular. Unlike \(\operatorname{ver}\), \(PAT\) shares with
\(P\) the property that weaker answers score higher (part~3)
\cite[p.~347]{Niiniluoto2011HintikkaProgramReading}.

\begin{hbexample}{def:p15-w02b-p15-030-nn-pat-worked}{Positive probable approximate truth at zero posterior}
Let \(B=\{h_1,h_2,h_3\}\) with \(\Delta_{12}=\Delta_{21}=0.1\) and all other
off-diagonal distances \(1\); take \(\gamma=1\), \(\gamma'=1/2\), \(g=h_1\),
\(\varepsilon=0.2\), and evidence with \(P(h_1/e)=0\), \(P(h_2/e)=0.7\),
\(P(h_3/e)=0.3\). Then \(P(g/e)=0\). Since \(\Delta_{\min}(h_2,g)=0.1\leq0.2\)
and \(\Delta_{\min}(h_3,g)=1\), \(V_{0.2}(g)=\{h_1,h_2\}\) and
\(PAT_{0.8}(g/e)=0.7\). By part~2 of
\cref{thm:p15-w02b-p15-030-truthlikeness}, with \(h_2\) as target the sum of
distances is \(1.1\), so \(Tr(h_1,h_2)=1-(1+0.5/1.1)\cdot0.1=47/55\approx0.855\);
with \(h_3\) as target the sum is \(2\), so \(Tr(h_1,h_3)=1-(1+0.25)\cdot1=-1/4\);
and \(Tr(h_1,h_1)=1\). Hence the probability that \(g\) is truthlike at least
to degree \(0.8\) is \(P(h_2/e)=0.7\), while
\(\operatorname{ver}(g/e)=0.7\cdot47/55-0.3/4=1151/2200\approx0.523\).
The refuted complete answer \(h_1\) thus keeps substantial probable
approximate truth and probable verisimilitude because its near neighbour
\(h_2\) is probable.
\end{hbexample}

Probable approximate truth, Niiniluoto argues, is essentially the PAC
notion of machine learning (Niiniluoto 2005). The AI community's
``concepts'' correspond to monadic constituents, so concept learning can
be modelled by Hintikka-style generalization; AI concept learning is
purely eliminative (positive and negative instances), whereas Hintikka's
system adds enumerative and analogical considerations. On formal learning
theory: (13) only says that degrees of belief in \(C^c\) converge to
certainty, not that \(C^c\) is the true \(C^*\), and (73) does not make
\(Tr(C^c,h^*)\) maximal; the stronger results need ESC (Niiniluoto 1987,
p.~276). Formal learning theory's convergence results are no stronger,
despite demanding success on all data streams, because they presuppose
something like ESC: decidability in the limit assumes data streams
``complete in that they exhaust the relevant evidence'' (Earman 1992,
p.~210) or ``perfect'' (Kelly 1996, p.~270, as cited); without ESC even
global underdetermination cannot be avoided (Kelly, p.~17), since an
infinite stream of white swans may be drawn from a world with black swans
too \cite[pp.~347--348]{Niiniluoto2011HintikkaProgramReading}.

More formally, Hintikka's learner uses epistemic priors \(P(C^w)\) and
likelihoods \(P(e/C^w)\), and (13) needs no assumption of objective sampling
conditions; the same holds for de Finetti's and Savage's results on the
convergence of opinion among agents starting from different
non-dogmatic priors (Howson and Urbach 1989; Earman 1992; compare
Niiniluoto 1984, p.~102). A footnote rejects Kuipers's (1977) restriction of the
\(K\)-dimensional system to multinomial contexts with objective
distributions, while allowing that such a condition may be a contextual
presupposition. One may add a fair sampling assumption (Kuipers 1977)
giving each kind a nonzero objective chance; such propensities do not
satisfy Plenitude and do not exclude sequences violating ESC, but these are
extremely improbable by the convergence theorems of the probability
calculus (compare Festa 1993, p.~76). For fair draws with
replacement from an urn \(w\) with proportion \(r\) of \(A\), the Strong Law of
Large Numbers gives \(k/n\to r\) with probability one, which Niiniluoto calls
weaker than ordinary convergence, because non-typical sequences violating
ESC cannot be logically excluded though their measure is zero. His
conclusion: formal learning theory and probabilistic induction are in the
same boat; ESC is why induction is fallible even in the ideal limit, and a
fallibilist ``convergent realist'' should claim neither decidability in the
limit nor gradual decidability but convergence to the truth with
probability one \cite[p.~348 and n.~11]{Niiniluoto2011HintikkaProgramReading}.
These are methodological claims; parts~4--6 of
\cref{thm:p15-w02b-p15-030-pat} is the only formal content proved here,
and the Strong Law is used as Niiniluoto cites it, not reproved.

\subsection{Niiniluoto's evaluation of the program}
\label{sec:p15-w02b-p15-030-nn-evaluation}

Inductive systems form a sequence of successive ``theories'' of rational
inference with notable continuity: the \(K\)-dimensional system generalizes
Hintikka's generalized combined system, which generalizes and modifies
Carnap's \(c^*\), which generalizes Laplace's rule; Hacking (1975, p.~141)
suggested a Leibnizian ought to like \(c^*\). Each successor contains its
predecessors as special or limiting cases and solves more problems, so the
sequence is progressive by Lakatos's standards: Hintikka's system handles
generalization, the \(K\)-dimensional system explains why Laplace-to-Carnap
systems failed at it, and the framework extends to analogy, observational
error and theoretical premises and conclusions
\cite[p.~349]{Niiniluoto2011HintikkaProgramReading}.

Hintikka's philosophical conclusion was that inductive probabilities are
always relative to extra-logical parameters: small \(\alpha\) makes the
posterior of generalizations grow fast and indexes boldness, large \(\alpha\)
caution; alternatively \(\alpha\) expresses a regularity assumption about the
lawlikeness of \(U\) (relative, a footnote adds, to the expressive power of
universal generalizations; Walk 1966 argues Carnap's \(\lambda\) is directly
proportional to the statistical entropy of the universe). This, for
Hintikka, ends hopes for purely logical probabilities in the sense of C3,
C5 and C6 (Hintikka 1987, 1997, p.~317). After the \(\lambda\)-continuum
Carnap's position was in fact similar: for processes with objective
probabilities there is an optimal \(\lambda\) (1952), whose estimation Festa
(1995) elaborates; in his later work (1980) Carnap tied \(\lambda\) to
attribute distance yet admitted it may reflect the investigator's
personality, concluded that \(\lambda\) should lie between \(\tfrac12\) and
\(K\), and finally recommended \(\lambda=1\)
\cite[p.~349 and n.~12]{Niiniluoto2011HintikkaProgramReading}.

The extra-logical character brings inductive logic close to Bayesianism:
it is Bayesianism that studies epistemic probabilities in relation to
structural assumptions (expressive power, parameters, order and disorder,
similarity), which in Suppes's terms (1969, p.~95) are not part of the pure
theory of rationality but limit where the theory applies. It finds more
structure than personalism and, rather than one system or one objective
prior per problem, offers a rich framework for many methodological
situations \cite[pp.~349--350]{Niiniluoto2011HintikkaProgramReading}.

\begin{hbexample}{def:p15-w02b-p15-030-nn-parameter-space}{Constituents versus statistical hypotheses}
A statistician might replace the constituent ``all individuals are \(Q_1\)''
by the hypothesis that the proportion of \(Q_1\) lies in a short interval
\((q,1]\) and use Bayesian statistics. But one counterexample refutes the
constituent, not its statistical counterpart; and Hintikkian priors put
positive weight on sets of geometric measure zero. With two
\(Q\)-predicates the parameter space is \(\{\langle p_1,p_2\rangle\mid
p_1+p_2=1\}\), and the three consistent constituents are
\[
 C_1=\{\langle1,0\rangle\},\qquad C_2=\{\langle0,1\rangle\},\qquad
 C_3=\{\langle p,1-p\rangle\mid0<p<1\}.
\]
Laplace's uniform prior gives \(C_1\) and \(C_2\) probability zero, since
each is a single point, and so \(C_3\) probability one; Hintikka's priors
give the two singletons positive mass. This illustrates how Hintikka's
improvement of Carnap breaks with standard Bayesian statistics
\cite[p.~350]{Niiniluoto2011HintikkaProgramReading}.
\end{hbexample}

Hintikka (1997) adds that inductive logic with parameters is not purely
Bayesian, since changing \(\lambda\) and \(\alpha\) is not conditionalization,
so inductive logic can study language change (compare Niiniluoto and
Tuomela 1973); and he proposed making all
extra-logical assumptions explicit premises, turning inductive logic into
part of deductive logic, a view repeated in his ``Replies'' (Auxier and
Hahn 2006). Later Hintikka became a critic of traditional treatments of
induction, developing an interrogative model of inquiry as a dialogue
between the inquirer and ``Nature''; there induction plays a modest role at
the micro-level of experiment, Hume's problem seems to vanish from serious
theory of method, Nature's answers in controlled experiments are already
general (at least of AE-complexity), and the ``Atomistic Postulate''
restricting answers to atomic propositions is to be rejected
\cite[pp.~350--351]{Niiniluoto2011HintikkaProgramReading}.

Niiniluoto finds this study of inquiry fruitful but doubts that inductive
generalization can be dispensed with everywhere by explicit strong
regularity assumptions. A footnote quotes Hintikka's reply (Auxier and Hahn
2006, p.~778) that once these premises are explicit no independent rules
of inductive inference are needed, and objects that this ignores the
distinction between factual and probabilistic detachment (note~8).
Hintikka's own system is not tied to the Atomistic Postulate as Carnap's
is, since it handles \(P(g/e\&T)\) and \(P(g/T)\) with a background theory
\(T\) \cite[p.~351 and n.~13]{Niiniluoto2011HintikkaProgramReading}.

His closing assessment: Carnap's logical-empiricist motivation looks
outdated, but Hintikka's handling of real generalizations and theories
changes the picture; inductive logic can be ``non-inductivist'' (Niiniluoto
and Tuomela 1973), and with its applications to theoretical
systematization, semantic information, explanation, abduction and
truthlikeness it can serve a critical scientific realist. Against Lakatos's
mockery of exact numerical degrees of confirmation, quantitative measures
serve mainly comparative judgments, and inductive logic studies general
principles of probabilistic and uncertain reasoning rather than supplying
calculation manuals (Niiniluoto 1980, p.~226). Although Hintikka's system is
not yet fully extended to first-order languages, distributive normal forms
suggest that the main features of monadic induction carry over. Monadic
cases can be implemented as programs; with computer science, robotics and
AI, inductive logic may be useful for machine learning (a footnote recalls
Carnap's 1971 ``inductive robots''), the AI community is reinventing logical
approaches proposed by philosophers, and the future paradigm of induction
may be found in machine learning
\cite[pp.~351--352 and n.~14]{Niiniluoto2011HintikkaProgramReading}.

\begin{proofboundary}{def:p15-w02b-p15-030-nn-boundary}{What the Niiniluoto reconstruction establishes}
Proved here, under the hypotheses stated in each theorem: the elementary
properties of Carnap's representative function; von Wright's principal
theorem (7); the derivation of (16), (16\('\)) and the posterior formulas
(20), (22), (23) with the asymptotics (13), (14) for the generalized
combined system; (18), and (19) and (25) for the priors (17\('\)); the
Kemeny--Carnap characterization on singular sentences; consequences of the
stated \(K\)-dimensional formulas and the intersection condition with the
two-dimensional system; the information identities; positive relevance
and the monotonicity (41)--(42) with the hypothesis analysis; the
expected-utility identities; the analogy formulas (58)--(60); and the
properties of \(Tr\), \(\operatorname{ver}\) and \(PAT\), including (66), (71)--(73)
and their ESC versions. Reported from the chapter without proof: the
Hintikka--Niiniluoto \(K\)-dimensional representation (the parameters,
(26), the product formula for \(P(C^K)\)), (27), (28) and the derivation of
(3) and (4) from (27); (19) and (25) for general \(\lambda\) and the
\(\alpha\to\infty\) limit of the whole two-dimensional continuum; the
approximations (21), (24), (56) and the claims depending on them; (31); the
threshold formula for \(n_0\); (43); and the results of Hilpinen, Tuomela,
Uchii, Kuipers and Niiniluoto--Tuomela cited in the chapter. Two printed
formulas are flagged rather than silently changed: the special values of
\(\lambda\) on p.~316, and the \(K\)-dimensional parameter formula on p.~324,
which conflicts with (27); and the printed (54)--(55) pair is consistent
only for \(q=1\).
\end{proofboundary}
